\begin{savequote}
\includegraphics[width=5cm]{./images/nasa-pd/154960main_image_feature_638_ys_full-00800.jpg}

Bild: NASA
\end{savequote}
%%%%%%%%%%%%%%%%%%%%%%%%%%%%%%%%%%%%%%%%%%%%%%%%%%%%%%%%%%%
%%%%%%%%%%%%%%%%%%%%%%%%%%%%%%%%%%%%%%%%%%%%%%%%%%%%%%%%%%%
%%%%%%%%%%%%%%%%%%%%%%%%%%%%%%%%%%%%%%%%%%%%%%%%%%%%%%%%%%%
%%%%%%%%%%%%%%%%%%%%%%%%%%%%%%%%%%%%%%%%%%%%%%%%%%%%%%%%%%%
%%%%%%%%%%%%%%%%%%%%%%%%%%%%%%%%%%%%%%%%%%%%%%%%%%%%%%%%%%%
\chapter
[\CO{[VIT]} Vitalfunktionen]
{\CC{[VIT]}\\Vitalfunktionen}
\label{chp:VIT}
% \AasRsN{RS~V}


% \emph{Maintainer: Sebastian Gabriel}

% \AuthNotes{
% Changelog:
% 
% \begin{description}
%     \item[2009-02-25] Start Changelog. Kapitel mit Korr
%     Koch korr. 
%     \item[2009-mm-dd] in \LaTeX übernommen
%     \item[2009-04-19] anhand der OpenOffice-Version überarbeitet.
% Bilder fehlen noch. 
% \item[2009-05-06] GAB: Abschnitte "`Blutdruck"' und
% "`Atemvolumina"' aus Anatomie-Teil übernommen.
% \end{description}
% 
% }

\ChapterIntro
{Diese Kapitel behandelt die Vitalfunktionen
des menschlichen Körpers. Sie ermöglichen die
Funktion des Körpers. Es gibt grundlegende Vitalfunktionen 1. Ordnung (Bewusstsein,
Atmung, Kreislauf), deren Ausfall binnen kurzem zum
Tod führen können, und Vitalfunktionen 2. Ordnung, deren
Störung längere Zeit toleriert werden kann.

Die Beurteilung und Erhaltung von Vitalfunktionen ist eine
wesentliche Aufgabe in der Sanitätshilfe.
}{%
\begin{description}
  \item[Maintainer:] S. Gabriel
  \item[Autoren:] Diverse
  \item[Reviewer:] Standard-Reviewprozess
  \item[Version:] {\DocVersion} ({\DocVersionInfo})
  \item[SHA1:] {\DocGitCommit}
\end{description}

{\TxTeilkapitelAASS}
}



\clearpage

\section{Einleitung}
\label{topic:vitalfunktionen-definition}
\label{topic:naca}
% Kapitel Einleitung hinzugeführt, um label sinnvoll setzen
% zu können (Emhofer)



\paragraph{Überblick}

Die Vitalfunktionen ermöglichen die
Funktion des Körpers. Man unterscheidet grundlegende
Vitalfunktionen 1. Ordnung (Bewusstsein, Atmung, Kreislauf), deren Ausfall binnen kurzem zum
Tod führen können, und Vitalfunktionen 2. Ordnung, deren
Störung längere Zeit toleriert werden kann:

\begin{table}[H]
\FormatTable
\Captionof{table}{Übersicht über die Vitalfunktionen 1. und
2. Ordnung}

\begin{tabularx}{\linewidth}[H]{XX}
\TabHeading{Vitalfunktionen I. Ordnung}
 &
\TabHeading{Vitalfunktionen II. Ordnung}
\\\addlinespace\midrule\addlinespace
\begin{itemize}
  \item Bewusstsein
  \item Atmung
  \item Kreislauf
\end{itemize}
\EE{Wichtigste, grundlegende Funktionen des Körpers

Ausfall einer dieser Funktionen bedeutet immer akute
Lebensgefahr und muss dringlich behandelt werden!}
 &
 \begin{itemize}
   \item Wasser-/Elektrolythaushalt
   \item Säure-/Basenhaushalt
   \item Temperaturhaushalt
   \item Hormonsystem
   \item Stoffwechsel
   \item Immunsystem
 \end{itemize}
\\\addlinespace\bottomrule
\end{tabularx}
\end{table}


\Layout{2}{NACA-Score}
{\index{NACA-Score}
\Z{\Lang{Abkz.} National Advisory Committee for Aeronautics

Der NACA-Score diente ursprünglich der Bewertung von 
Erkrankungen / Verletzungen als Entscheidungshilfe
welche Soldaten zuerst auszufliegen wären.} Heutzutage dient
der NACA-Score der groben statistischen Erfassung.

Man unterscheidet die Stufen 1 bis 7:


\begin{table}[H]
\FormatTable
\Captionof{table}{NACA-Score.}

\begin{tabular}{clrr}\addlinespace
\TabHeading{NACA}
&
\TabHeading{Bedeutung}
&
\TabHeading{Anteil Pat.}
&
\TabHeading{Anteil \PC}
\\\addlinespace\midrule\addlinespace
 I & keine ärztliche Therapie erforderlich & 10\,217 & 8,7
\\\addlinespace
 II & ambulante Abklärung & 62\,207 & 52,8
\\\addlinespace
 III & stationäre Abklärung & 38\,307 & 32,5
\\\addlinespace
 IV & akute Lebensgefahr nicht ausgeschlossen & 2\,830 & 2,4
\\\addlinespace
 V & akute Lebensgefahr & 1\,022 & 0,9
\\\addlinespace
 VI & Reanimation & 34 & 0,03	
\\\addlinespace
 VII & Tod & 312 & 0,3
 \\\addlinespace\midrule\addlinespace
 &&117\,849 & 100
\\\addlinespace\bottomrule
\end{tabular}

\Z{Die angegebenen Fallzahlen beziehen sich auf Wien im Jahr
2009, nach  \citep{KontrollamtWienKA-K-13-09}.}
\end{table}
}{}{}{}{}



\clearpage
%%%%%%%%%%%%%%%%%%%%%%%%%%%%%%%%%%%%%%%%%%%%%%%%%%%%%%%%%%%
%%%%%%%%%%%%%%%%%%%%%%%%%%%%%%%%%%%%%%%%%%%%%%%%%%%%%%%%%%%
%%%%%%%%%%%%%%%%%%%%%%%%%%%%%%%%%%%%%%%%%%%%%%%%%%%%%%%%%%%
%%%%%%%%%%%%%%%%%%%%%%%%%%%%%%%%%%%%%%%%%%%%%%%%%%%%%%%%%%%
\section{Bewusstsein}
\label{topic:bewusstsein}
\index{Bewusstsein}

\Definition{Oberbegriff für u.\,a. Wachheit,
Orientierung, Aufmerksamkeit, Auffassungsgabe, Denkverlauf und
Merkfähigkeit \cite{Pschy259}.}

\Layout{abstract}{}
{Das Bewusstsein ist eine grundlegende Vitalfunktion erster
Ordnung. Bei der Diagnose beurteilt man die
Bewusstseinsgrade (klar, somnolent, soporör, komatös) und
die Orientierung (zeitlich, örtlich, zur Situation und zur
Person). Bewusstseinsstörungen sind aufgrund des Ausfalls
der Schutzreflexe und des Zurückfallens der Zunge
gefährlich. Es gibt Ursachen für Störungen, welche primär
direkt (Verletzung, Schlaganfall, \ldots), und Ursachen, die sekundär, also über Umwege (Zuckerstoffwechsel, Vergiftungen, \ldots), auf das ZNS einwirken.

}{}{}{}{}

\Z{Die "`wahre"' Definition von Bewusstsein ist seit
tausenden von Jahren ungeklärt. Auch an dieser Stelle wird
das Geheimnis nicht gelüftet werden, daher muss die obige
Definition für den Alltag ausreichen.\footnote{Hinweisen
zufolge könnte die Antwort
eventuell \enquote{42} sein
\citep{UltimateHitchhikersGuide96hitchhikersguide}.}}

% \protect\DefinitionExp{Bewusstsein}
% 
% bla
% 
% \protect\glsentrydesc{Bewusstsein}

Das Bewusstsein ist besonders wichtig hinsichtlich des
Schutzes des Menschen vor Bedrohung. Ein bewusstseinsklarer
Mensch kann sich gegen Gefahren wehren, ein eingetrübter
oder bewusstloser Mensch kann dies schlecht oder gar nicht.

%%%%%%%%%%%%%%%%%%%%%%%%%%%%%%%%%%%%%%%%%%%%%%%%%%%%%%%%%%%
%%%%%%%%%%%%%%%%%%%%%%%%%%%%%%%%%%%%%%%%%%%%%%%%%%%%%%%%%%%
%%%%%%%%%%%%%%%%%%%%%%%%%%%%%%%%%%%%%%%%%%%%%%%%%%%%%%%%%%%
\subsection{Bewusstseinsstörungen}
\label{topic:bewusstseinsstoerung}

\Definition{Störung des Bewusstseins, besonders der
Wachheit und Orientierung.}


%%%%%%%%%%%%%%%%%%%%%%%%%%%%%%%%%%%%%%%%%%%%%%%%%%%%%%%%%%%
\Layout{0}{Allgemein}
{%
Wie bereits oben erwähnt, ist das Bewusstsein wesentlich an
der Abwehr von Gefahren beteiligt. Je mehr das Bewusstsein
gestört oder reduziert ist, desto gefährlicher ist das für
den Patienten.

Die \EE{Zunge} kann zurück fallen und die Atemwege verlegen.
Dabei ist manchmal ein dem \E{Schnarchen ähnliches}
Atemgeräusch zu vernehmen.

Durch den Ausfall \index{Schutzreflexe!Ausfall} der
Schutzreflexe kann es zu einer
\II{Aspiration} von Speichel, Blut oder
Erbrochenen kommen. 
}{%
Gefahren:
\begin{itemize}
  \item Zurückfallen der Zunge
  \item Ausfall der Schutzreflexe
  \item Aspiration
\end{itemize}
}{}{}{}

%%%%%%%%%%%%%%%%%%%%%%%%%%%%%%%%%%%%%%%%%%%%%%%%%%%%%%%%%%%
\Layout{2}{Schutzreflexe}
{%
Ab einem gewissen Punkt können auch so wesentliche
Funktionen wie der \EE{Schluck-}, \EE{Würge-} bzw. \EE{Hustenreflex}
ausfallen, und es kann Mageninhalt oder Blut in die
Lunge gelangen (\T{Aspiration}). Man spricht auch vom
\EE{Ausfall der Schutzreflexe}.

Bei schweren Bewusstseinsstörungen kann es auch zum Erschlaffen
wichtiger Muskeln wie z. B. der Zunge kommen. Fällt
diese zurück, kann sie den Atemweg blockieren oder ganz verschließen.
}{}{}{}{}


\Fazit{Schwere Bewusstseinsstörungen sind
\EE{lebensgefährlich!}}


%%%%%%%%%%%%%%%%%%%%%%%%%%%%%%%%%%%%%%%%%%%%%%%%%%%%%%%%%%%
\Layout{0}{Einteilung der Störungen}
{%
Man unterscheidet zwischen der Quantität ("`\E{Wieviel}
Bewusstsein ist vorhanden?"') und der Qualität ("`\E{Wie
gut} funktioniert das vorhandene Bewusstsein?"'). Maßgeblich
ist der \EE{Bewusstseinsgrad} (quanititativ) und die
\EE{Orientiertheit} des Patienten (qualitativ).

Die Glasgow Coma Scale ist ein spezielles Maßsystem für
vorwiegend statistische Zwecke.
}{
\begin{itemize}
  \item Bewusstseinsgrad (quantitativ)
  \item Orientiertheit (qualitativ)
\end{itemize}
}{}{}{}

%%%%%%%%%%%%%%%%%%%%%%%%%%%%%%%%%%%%%%%%%%%%%%%%%%%%%%%%%%%
\Layout{2}{Bewusstseinsgrad}
{%
\index{Bewusstseinsgrad}
\label{topic:bewusstseinsstoerung-grade}
\index{Bewusstsein!Symptome}
%
Der Bewusstseinsgrad gibt Auskunft über die Quantität des Bewusstseins ("`\E{Wieviel} Bewusstsein ist
vorhanden?"'). Er muss bei jedem Patienten festgestellt werden.

Die folgende Tabelle (\prettyref{tab:bewusstseinsgrade}) listet die möglichen
Bewusstseinsgrade auf:
}{
\begin{itemize}
  \item Klar
  \item Getrübt
  \begin{itemize}
    \item Somnolent
    \item Soporös
  \end{itemize}
  \item Bewusstlos -- Koma
\end{itemize}
}{}{}{}

\begin{table}[H]
\FormatTable
\Captionof{table}{Bewusstseinsgrade\label{tab:bewusstseinsgrade}}

\begin{tabulary}{\linewidth}{LlLL}
% \toprule
\addlinespace\toprule\addlinespace
\TabSub{Bewusstseinsklar}
 &
 &
 wach
 &
\\\addlinespace
\TabSub{Bewusstseinsgetrübt}
&
\noindent\T[Somnolenz]{somnolent}
&
schläfrig, aber leicht erweckbar.
 &
\\
 &
\noindent\T[Sopor]{soporös}
 &
 kaum und nur mit erheblichen Aufwand (Schmerzreiz)
 erweckbar 
 &
\noindent\EE{Gefahr!}
\\\addlinespace
\TabSub{Bewusstlos}\index{Bewusstlosigkeit}
 &
\noindent\T[Koma]{komatös}
 &
 nicht erweckbar
 &
\noindent\Ca{Lebensgefahr!}
\\\addlinespace\bottomrule
\end{tabulary}
\end{table}

%%%%%%%%%%%%%%%%%%%%%%%%%%%%%%%%%%%%%%%%%%%%%%%%%%%%%%%%%%%
\Layout{2}{Orientierung}
{%
\label{topic:orientierung}
\index{Orientierung}
%
Es ist wichtig zu Überprüfen ob der Patient orientiert ist.
Man unterscheidet dabei die Orientierung zu:  

\begin{itemize}
  \item \EE{Person}: Name, Alter.
  \item \EE{Zeit}: Tageszeit,
  Wochentag, Monat, Jahreszeit, Jahr?
  \item \EE{Ort}: Wo befindet sich der Patient gerade?
  \newline
  (Stadt, Land, Adresse, Wohnung, Auto, etc.)
  \item \EE{Situation}: Was passiert gerade?
\end{itemize}

}{
\begin{itemize}
    \item Zur Person
    \item Zeitlich
    \item Örtlich
    \item Zur Situation
\end{itemize}
}{}{}{}

%%%%%%%%%%%%%%%%%%%%%%%%%%%%%%%%%%%%%%%%%%%%%%%%%%%%%%%%%%%
\Layout{2}{Glasgow Coma Scale (GCS)}
{%
\index{Glasgow Coma Scale}
\index{GCS}
\label{topic:gcs}
%
Die Glasgow Coma Scale ist eine spezielle Skala, welche der
'objektiven' Einschätzung der Bewusstseinslage  dient. Sie
hat einen Wertebereich von \VonBis{3}{15} Punkten. Ein voll
bewusstseinsklarer Patient hat 15 Punkte, ein schwerst
bewusstseinsgestörter oder toter Patient 3 Punkte (nicht
0!). Bewertet werden die Augenöffnung, die beste verbale und
die beste motorische Reaktion des Patienten.

Heutzutage dient die \Abk{GCS} vorwiegend statistischen
Zwecken\ZFootnote{Regelmäßig findet man Handlungsanweisungen
-- überwiegend für ärztliches Personal -- in der Literatur,
welche die \Abk{GCS} als Entscheidungskriterium (z.\,B. für
die Intubation) heranziehen. Der Nutzen ist aber genauso
regelmäßig umstritten.}.
}{
\VOnBis{3}{15} Punkte.
\begin{itemize}
    \item Augenöffnung
    \item Beste verbale Reaktion
    \item Beste motorische Reaktion
\end{itemize}
\prettyref{tab:gcs}.
}{}{}{}


\begin{table}[H]
\FormatTable
\Captionof{table}{GCS - Glasgow Coma Scale für den Erwachsenen.\label{tab:gcs}}

\begin{tabulary}{\linewidth}{LCCL}\addlinespace\toprule\addlinespace
\noindent ~ & \noindent\TabHeading{Max.} & \noindent\TabHeading{Punkte} &
\noindent\TabHeading{Reaktion}
\\\midrule
\TabSub{Augenöffnung} & 4 & 4 & spontan
\\
 & & 3 & auf Ansprache
\\
 & & 2 & Schmerzreiz
\\
 & & 1 & nicht
\\\midrule
\TabSub{Verbale Antwort} & 5 & 5 & klar
\\
 & & 4 & verwaschen
\\
 & & 3 & stammelnd
\\
 & & 2 & unartikuliert
\\
 & & 1 & nicht
\\\midrule
\TabSub{Motorische Reaktion} & 6 & 6 & spontan
\\
 & & 5 & gezielte Abwehrbewegungen
\\
 & & 4 & ungezielte Abwehrbewegungen
\\
 & & 3 & Beugung
\\
 & & 2 & Strecken
\\
 & & 1  & bewegt gar nicht
\\\bottomrule
\end{tabulary}
\end{table}




%%%%%%%%%%%%%%%%%%%%%%%%%%%%%%%%%%%%%%%%%%%%%%%%%%%%%%%%%%%
\Layout{0}{Ursachen}
{%
Durch die Schädigung des zentralen Nervensystems kommt es zu
der Bewusstseinsstörung. Diese Schädigung kann direkt das
\Abk{ZNS} betreffen (primäre Ursache) oder über Umwege einen
schädlichen Einfluss auf das \Abk{ZNS} haben (sekundäre
Ursache). \prettyref{tab:bewusstseinsstoerungen-ursachen}
listet typische Ursachen auf.

\begin{table}[H]
\FormatTable
\Captionof{table}{Typische Ursachen von
Bewusstseinsstörungen.\label{tab:bewusstseinsstoerungen-ursachen}}

\begin{tabularx}{\linewidth}[H]{XX}\addlinespace\toprule\addlinespace
\noindent\TabHeading{Primäre Ursachen} & \noindent\TabHeading{Sekundäre
Ursachen}
\\
{\tiny betreffen direkt das ZNS} & {\tiny betreffen über
Umwege das ZNS} 
\\\addlinespace\cmidrule(lr){1-1}\cmidrule(lr){2-2}\addlinespace
\begin{itemize}
    \item Verletzungen (SHT)
    \item Blutungen, Hirnschwellung
    \item Schlaganfall
    \item Epileptischer Anfall
    \item Entzündungen
    \item \Z{Tumor}
    \item \ldots
\end{itemize}
 &
 \begin{itemize}
   \item Störungen der \ce{O2}-Versorgung:
   \begin{itemize}
     \item Atemstörungen  (Hypoxie)
     \item Herz-Kreislaufstörungen
   \end{itemize}
   \item Störungen des Wasserhaushaltes und
   Elektrolytentgleisungen
   \item Stoffwechselstörungen
   \begin{itemize}
     \item Zuckerhaushalt
     \item Leberversagen
   \end{itemize}
   \item Vergiftungen:
   \begin{itemize}
     \item Alkohol
     \item Medikamente, Drogen
     \item Reizgase,  Lösungsmittel ...
   \end{itemize}
   \item \ldots
 \end{itemize}
\\\bottomrule
\end{tabularx}
\end{table}
}
{
\begin{itemize}
  \item Direkt das \Abk{ZNS} betreffend (primär)
  \item Indirekt das \Abk{ZNS} betreffend (sekundär)
\end{itemize}

\prettyref{tab:bewusstseinsstoerungen-ursachen}.
}{}{}{}










\MassnahmenMK{Spezielle Maßnahmen}{Bewusstseinseintrübung}
{m:bewusstseinstruebung}{}{}

\MassnahmenMK{Spezielle Maßnahmen}{Bewusstlose und
soporöse Patienten}{m:bewusstlosigkeit}
{%
}{}
% \MassnahmenNG{Spezielle Maßnahmen}{Bewusstseinseintrübung}
% {\label{m:bewusstseinstruebung}}{ 
% % \paragraph{Maßnahmen
% % bei getrübten Patienten}~
% 
% \begin{itemize}
%   \item Lagerung: situationsgerecht je nach
%   Verdachtsdiagnose, im Zweifel \EE{stabile
%   Seitenlage}, bei hochschwangeren Frauen auf die
%   \E{linke} Seite.
%   \item \TxBeiVit \TxMassVitBes
%   
%   \begin{itemize}
%     \item \ce{O2}: nach Verdachtsdiagnose, im Zweifel
%     ca. \VonBis{6}{8}\LPerMin*
%     \item \EE{Monitoring}: Auf \EE{Bewußtseinsgrad}
%     besonders achten!
%   \end{itemize}
%   
%   
%   \item Diagnostische Schwerpunkte:
%   \begin{itemize}
%     \item \EE{Ursachenforschung} -- Warum ist der
%     Patient eingetrübt?
%     \item (Fremd-)Anamnese, Umstände, Szenerie
%     \item \EE{Neurocheck}, inkl. \EE{BZ}-Messung!
%     \item \EE{Traumacheck}. Suche nach Primär- und
%     Folgeverletzungen!
%   \end{itemize}
% \end{itemize}
% 
% }
% 
% \MassnahmenNG{Spezielle Maßnahmen}{Bewusstlose und
% soporöse Patienten}{\label{m:bewusstlosigkeit}}
% {%
% \begin{itemize}
%   \item \TxMassVitBes
%   \begin{itemize}
%     \item Lagerung: stabile Seitenlage, bei hochschwangeren Frauen auf die
%     \E{linke} Seite.
%     \item \ce{O2}:  10\LPerMin* , bei
%     Zyanose\footnote{Zyanose: Bläuliche Verfärbung der Haut
%     infolge einer unzureichenden Sauerstoffversorgung.
%     Besonders auffällig an den Lippen, Gesicht, und
%     Extremitäten.}  15\LPerMin*
%     \item \EE{Monitoring}: Auf \EE{Bewußtseinsgrad}
%     besonders achten!
%   \end{itemize}
%   \item Diagnostische Schwerpunkte:
%   \begin{itemize}
%     \item \EE{Ursachenforschung} -- Warum ist der Patient bewusstlos oder
%     soporös?
%     \item (Fremd-)Anamnese, Umstände, Szenerie
%     \item \EE{Neurocheck}, inkl. \EE{BZ}-Messung!
%     \item \EE{Traumacheck}. Suche nach Primär- und
%     Folgeverletzungen!
%   \end{itemize}
% \end{itemize}
% }





\clearpage
%%%%%%%%%%%%%%%%%%%%%%%%%%%%%%%%%%%%%%%%%%%%%%%%%%%%%%%%%%%
%%%%%%%%%%%%%%%%%%%%%%%%%%%%%%%%%%%%%%%%%%%%%%%%%%%%%%%%%%%
%%%%%%%%%%%%%%%%%%%%%%%%%%%%%%%%%%%%%%%%%%%%%%%%%%%%%%%%%%%
%%%%%%%%%%%%%%%%%%%%%%%%%%%%%%%%%%%%%%%%%%%%%%%%%%%%%%%%%%%
\section{Atmung}

%%%%%%%%%%%%%%%%%%%%%%%%%%%%%%%%%%%%%%%%%%%%%%%%%%%%%%%%%%%
%%%%%%%%%%%%%%%%%%%%%%%%%%%%%%%%%%%%%%%%%%%%%%%%%%%%%%%%%%%
%%%%%%%%%%%%%%%%%%%%%%%%%%%%%%%%%%%%%%%%%%%%%%%%%%%%%%%%%%%
\subsection{Atemfrequenz, -tiefe und -volumen}
\label{topic:atemvolumina}


% STEU 2009-06-28: Trennung Oxygenation und Ventilation
\A{STEU 2009-06-28: Trennung Oxygenation und Ventilation}

\Layout{0}{Kennzahlen}
{%
\index{Atemfrequenz}
\index{AF}
\index{Atemzugsvolumen}
\index{AZV}
\index{Atemminutenvolumen}  
\index{AMV} 
%
Es gibt 3 elementare Kennzahlen, welche die wichtigsten Atmungsparameter
beschreiben:

\begin{labeling}[\dots{}]{Atemminutenvolumen (AMV)}
  \item[\T{Atemfrequenz} (\T{AF})] gibt die Anzahl der Atemzüge pro
  Minute an.
  \item[\T{Atemzugsvolumen} (\T{AZV})] gibt die Menge
  Luft an, die \E{pro Atemzug} eingeatmet wird.
  \item[\T{Atemminutenvolumen} (\T{AMV})] gibt die Menge an \E{pro Minute}
  eingeatmeter Luft an.
\end{labeling}

Das Atemminutenvolumen kann aus der Atemfrequenz und dem
Atemzugsvolumen errechnet werden:

\begin{equation}
\mbox{AMV} = \mbox{AZV} \times \mbox{AF}
\end{equation}

% Die
% \T{Atemfrequenz} (\T{AF}) gibt die Anzahl der Atemzüge pro
% Minute an. Das \T{Atemzugsvolumen} (\T{AZV}) gibt die Menge
% an Luft an die eingeatmet wird. Das \T{Atemminutenvolumen}
% (\T{AMV}) errechnet sich aus der Atemfrequenz und dem
% Atemzugsvolumen. Es gibt die Menge an \E{pro Minute}
% eingeatmeter Luft an.
}{
\begin{description}
    \item[AF] {\dotfill} Atemfrequenz
    \item[AZV] {\dotfill} Atemzugsvolumen
    \item[AMV] {\dotfill} Atemminutenvolumen
\end{description}

$\mbox{AMV} = \mbox{AZV} \times \mbox{AF}$
}{}{}{}



\Layout{0}{Normalwerte Atmung}
{%
Die normale Atemfrequenz (\T{AF}) beträgt beim Erwachsenen
in Ruhe \VonBis{12}{16}\PerMin*, das Atemzugsvolumen
(\T{AZV}) beträgt ca. 500\Ml*.\ZFootnote{Die Einstellungen eines Beatmungsgeräts können von den
Normalwerten deutlich abweichen, da man z.\,B. manchmal
"`vorsichtiger"', bzw. Lungen-schonender, beatmen möchte.} Das entspricht einem
Atemminutenvolumen (\T{AMV}) von  6\,--\,8\LPerMin*.  Für
eine ausreichende Atmung ist auch eine ausreichende
\ce{O2}-Konzentration der Umgebungsluft notwendig!
}{
\begin{description}
    \item[AF] {\dotfill} \EE{\VonBis{12}{16}\PerMin*}
    \item[AZV] {\dotfill} \EE{500\Ml*}
    \item[AMV] {\dotfill} \EE{\VonBis{6}{8}\Liter*}
\end{description}
}{}{}{}

Auch wichtig ist die ausreichende
\ce{O2}-Konzentration in der Umgebungsluft. Ohne Sauerstoff in der Einatemluft
ist die beste Atmung nutzlos. Der normale \ce{O2}-Gehalt der Umgebungsluft
beträgt 21\PC*.
\marginlabel{\cite{KlinkePhysio3,SchmidtPhysio28,DRAn3,StriebelAn7}}

Die genauen Werte sind abhängig von den Faktoren Alter,
Geschlecht und Größe. 


\Fazit{AMV = AZV $\ast$ AF, beim Erwachsenen  ca.
\VonBis{6}{8}\Liter*}





Die Atemfrequenz ist stark altersabhängig:



\begin{table}[H]
\Captionof{table}{Atemfrequenz im Altersvergleich,
Richtwerte.}

\begin{tabular}{lrr}\addlinespace\toprule\addlinespace
 & \noindent\TabHeading{AF} & \noindent\TabHeading{AZV}
\\\addlinespace\midrule\addlinespace
\TabSub{Neugeborenes} & 40\PerMin* & \VonBis{20}{40}\Ml*
\\\addlinespace 
\TabSub{Kleinkind} & 25\PerMin* & \VonBis{100}{200}\Ml* 
\\\addlinespace
\TabSub{Schulkind} & 20\PerMin* & \VonBis{200}{400}\Ml* 
\\\addlinespace
\TabSub{Erwachsener} & \VonBis{12}{16}\PerMin*  & ca. 500\Ml*
\\\addlinespace\bottomrule
\end{tabular}
\end{table}

\Layout{0}{Totraum}
{%
\index{Totraum}%
Unter dem Totraum versteht man den Teil der Atemwege der
nicht am Gasaustausch teilnimmt. Dies ist im wesentlichen
der gesamte Weg von der Nasenspitze bis kurz vor den
Lungenbläschen (Alveolen). Die Luft muss diesen Totraum
passieren um in die Lungenbläschen zu gelangen! Der Totraum
umfasst beim Erwachsenen ca. 150\Ml*. \cite{KlinkePhysio3}

% \index{Totraumventilation}
\T[Totraum!-ventilation]{Totraumventilation}: Bei der
Totraumventilation ist das Atemzugsvolumen so klein dass die
"`frische"' Luft im Totraum bleibt und nicht in die
Lungenbläschen gelangt: Es kommt keine frische Luft in die
Lungenbläschen und folglich kein Sauerstoff ins Blut.

\label{topic:schnappatmung}
Bei der \index{Schnappatmung}\T{Schnappatmung} kann dies
passieren: Auf den ersten Blick sieht es so aus als würde
der Patient "`schnappend"' atmen, in Wirklichkeit ist die
Atmung so gering dass es nur zu einer Totraumventilation
kommt, d.\,h. die Atmung ist nicht effektiv.
}{
\begin{itemize}
    \item Totraum: Teil der Atemwege, der nicht am
    Gassaustausch teilnimmt
    \item Totraumventilation
    \item Schnappatmung = nicht ausreichend
\end{itemize}
}{}{}{}



\Cave{Schnappatmung = keine Atmung}




\Layout{0}{Steuerung der Atmung}
{%
Die Atmung wird durch das \EE{Atemzentrum im Hirnstamm}
gesteuert. Daneben gibt es Sensoren im Körper, die den
Gasgehalt des Blutes beobachten.

Beim gesunden Patienten ist der \EE{Atemantrieb} vom
\EE{\ce{CO2}-Gehalt} des Blutes abhängig. Der
Sauerstoffanteil spielt nur eine untergeordnete Rolle! Unter
\FullPrettyRef{topic:saeure-basen-haushalt}, wird
ausgeführt, dass das \ce{CO2} eine wichtige Rolle im
Säure-Basen-Haushalt des Körpers einnimmt.

Wenn bei \EE{bestimmten Krankheitsbildern} die
\E{\ce{CO2}-Konzentration ständig hoch} ist, \E{gewöhnt}
sich der Körper an diesen Zustand. Er beginnt dann den
\E{Atemantrieb direkt über den \ce{O2}}-Gehalt im Blut zu
regeln.

Dies kann \EE{Probleme} machen: Wird bei so einem Patienten
der \ce{O2}-Gehalt plötzlich künstlich erhöht, fehlt der
Atemantrieb und der Patient kann \EE{Bewusstseinsstörungen}
(\I[\ce{CO2}-Narkose@CO2-Narkose]{\ce{CO2}-Narkose}) und
einen \EE{Atemstillstand} erleiden.

Das wichtigste Krankheitsbild bei dem so etwa vorkommen kann,
ist die \Abk{COPD} (\ref{topic:copd}).
}{
\begin{itemize}
  \item Atemzentrum im Hirnstamm
  \item Atemantrieb normalerweise durch \EE{\ce{CO2}}
  \item Krankhaft:
  \begin{itemize}
    \item Wenn \ce{CO2}-Konzentration dauerhaft erhöht \Sign{→} Regelung durch
    \ce{O2}-Gehalt.
    \item Problem bei Sauerstoffgabe
    \item V.\,a. bei \Abk{COPD} (\ref{topic:copd})
  \end{itemize}
\end{itemize}
}{}{}{}

\Fazit{Normalerweise wird der Atemantrieb über den
\ce{CO2}-Gehalt des Blutes gesteuert.}




%%%%%%%%%%%%%%%%%%%%%%%%%%%%%%%%%%%%%%%%%%%%%%%%%%%%%%%%%%%
%%%%%%%%%%%%%%%%%%%%%%%%%%%%%%%%%%%%%%%%%%%%%%%%%%%%%%%%%%%
%%%%%%%%%%%%%%%%%%%%%%%%%%%%%%%%%%%%%%%%%%%%%%%%%%%%%%%%%%%
\subsection{Störungen der Atmung}
\label{topic:atemstörungen}
\label{topic:atemstillstand}

\Layout{0}{Einleitung}
{%
\index{Ateminsuffizienz}%
\index{Insuffizienz!Atem-}%
\index{Dyspnoe}%
\index{Apnoe}%
\index{Bradypnoe}%
\index{Tachypnoe}%
Störungen der Atmung können aus vielen Ursachen eintreten.
\prettyref{tab:atemstoerungen-ursachen} beschreibt die
wichtigsten Pathomechanismen und Störungen.

Der Begriff \T{Ateminsuffizienz} beschreibt allgemein eine
\EE{nicht ausreichende Atemleistung}, unabhängig von der
Ursache der Störung. Dies führt zu einem \ce{O2}-Mangel
(Hypoxie). Oft hat der Patient das Gefühl nicht genug Luft
zu bekommen und leidet an \EE{Atemnot} (\T{Dyspnoe}).

Der \EE{Atemstillstand} (\T{Apnoe}) ist die extremste Form
der Ateminsuffizenz. Störungen der Atemgeschwindigkeit
werden als \E{Bradypnoe} (zu langsam) und als \E{Tachypnoe}
(zu schnell) bezeichnet.
}{
\begin{itemize}
  \item Ateminsuffizienz \ldots{} nicht ausreichende Atmung
  \item Hypoxie \ldots{} \ce{O2}-Mangel
  \item Atemnot (Dyspnoe)
  \item Atemstillstand (Apnoe)
  \item Bradypnoe (zu langsam)
  \item Tachypnoe (zu schnell)
\end{itemize}
}{}{}{}


\Layout{0}{Gefahr der Ateminsuffizienz}
{%
Bei einer Ateminsuffizienz sinkt der \EE{\ce{O2}-Gehalt} im
Blut. Es kommt zu einer \EE{Hypoxie} und in weiterer Folge
zu einer Unterversorgung lebenswichtiger Organe wie Herz und
Hirn.

Doch auch die \EE{\ce{CO2}-Abgabe} ist gestört und es kommt
zu einer Anhäufung von \ce{CO2} im Blut. In weiterer Folge
kann es zu einer \E{atmungsbedingten Übersäuerung}
(\EE{Respiratorische Azidose},
\prettyref{topic:azidose-respiratorische}) kommen.
}{
\begin{itemize}
    \item \ce{O2}-Gehalt im Blut ↓ (Hypoxie)
    \item \ce{CO2}-Gehalt im Blut ↑
    \item Respiratorische Azidose (\prettyref{topic:azidose-respiratorische})
\end{itemize}
}{}{}{}

\AuthNotes{GABG: pathologische Atmungstypen inkl.
Zeichnung einbinden? Kussmaul, etc.}



\Layout{0}{Ursachen}
{%
\prettyref{tab:atemstoerungen-ursachen}.
}{\prettyref{tab:atemstoerungen-ursachen}.}{}{}{}


\begin{table}[H]
\begin{WidePar}
\FormatTable
\Captionof{table}{Ursachen für
Atemstörungen.\label{tab:atemstoerungen-ursachen}}

\begin{tabulary}{\linewidth}[H]{LL}\addlinespace\toprule\addlinespace
\noindent\TabHeading{Störung des/der}  &
\noindent\TabHeading{Vorkommen und Erkrankungen}
\\\addlinespace\midrule\addlinespace
\TabSub{\ce{O2}-Angebotes}
 &
Gärkeller (\ce{CO2}\,${\uparrow}$, \ce{O2}\,${\downarrow}$)

"`dünne Luft"' (z.\,B. Hochgebirge,

p\ce{O2}\,${\downarrow}$)

Defekte Gastherme (\ce{CO}\,${\uparrow}$,
\ce{O2}\,${\downarrow}$)

Ertrinken
\\\addlinespace
\TabSub{\ce{O2}-Diffusion}
 &
Lungenödem

Lungenentzündung

Lungenembolie
\\\addlinespace
\TabSub{Atemmechanik}
 &
Verlegung der oberen Atemwege (Zunge, Erbrochenes, Laryngospasmus,
         Glottisödem, Bolus)

Verlegung der unteren Atemwege (Bronchitis, Asthma,
Lungenödem)

Verminderung der Dehnbarkeit der Thoraxwand oder des Lungengewebes
     (Rippenfraktur, Pneumothorax, Pleuraerguß, Emphysem)
\\\addlinespace
\TabSub{Neuro\-muskulären Atemregulation}
 &
SHT (Schädel Hirn Trauma)

Schlaganfall (Insult)

Vergiftungen

Tumore

Hohe Querschnittslähmung (Rückenmark durchtrennt)

Muskelerkrankungen

Nervenerkrankung

Stoffwechselstörungen
\\\addlinespace
\TabSub{Sonstige Regulationsstörung}
 &
 Psychogen (z.\,B. Panikattacke)
\\\addlinespace\bottomrule
\end{tabulary}
\end{WidePar}
\end{table}


\Layout{0}{Symptome}
{%
Atemstörungen können viele verschiedene Symptome
verursachen. \prettyref{tab:atmung-symptome} listet
verschiedene Symptome auf. Neben technischen Geräten ist vor
allem \E{sehen},
\E{hören}, \E{fühlen} und ein \E{aufmerksamer Blick}
gefordert.

Die Abgrenzung der Atemstörung zur Kreislaufstörung ist oft
nicht einfach. Eine Störung des Kreislaufes hat meist auch
einen verminderten Sauerstoff\E{transport} zur Folge, daher
kommt es dabei häufig \E{auch} zur Atemnot u.\,ä. Symptomen.
}{\prettyref{tab:atmung-symptome}.}{}{}{}


\begin{table}[H]
\begin{WidePar}
\FormatTable
\Captionof{table}{Symptome: Atmung.\label{tab:atmung-symptome}}

\begin{tabulary}{\linewidth}[H]{LLLL}
\toprule\addlinespace
\noindent\TabHeading{Kriterium}
 &
\noindent\TabHeading{Befund}
 &
\noindent\TabHeading{Beschreibung}
 &
%\noindent\TabHeading{Anm.}
\\\addlinespace\midrule\addlinespace
\TabSub{Beschwerde}
 &
Atemnot (Dyspnoe) \index{Dyspnoe}
 &
Leitsymptom
 &
\\\addlinespace
\TabSub{Atemgeräusch}
 &
Brodeln \index{Atemgeräusch!brodelndes}
 &
Blubbern, klassisch für Lungenödem
\index{Atemgeräusch!blubberndes} &
\\\addlinespace
 &
Stridor \index{Stridor}
 &
 &
\\\addlinespace
 &
Brummen, Giemen
\index{Atemgeräusch!brummendes}\index{Atemgeräusch!giemendes}
& &
\\\addlinespace
 &
Rasselgeräusche \index{Rasselgeräusche}
 &
 &
\\\addlinespace
\TabSub{Frequenz}
 &
Beschleunigt
 &
Tachypnoe \index{Tachypnoe}
 &
\\\addlinespace
 &
Verlangsamt
 &
Bradypnoe \index{Bradypnoe}
 &
\\\addlinespace
\TabSub{Atemzugs\-volumen}
 &
Schnappatmung
 &
\FullPrettyRef{topic:schnappatmung}
 &
\\\addlinespace
 &
Flache Atmung
 &
 &
\\\addlinespace
 &
Tiefe Atmung
 &
Z.\,B. Kußmaul'sche Atmung
 &
\\\addlinespace
\TabSub{Hautfarbe}
 &
Blass
 &
Normal, Anämie, Blutverlust
 &
\\\addlinespace
 &
Rosig
 &
Normal, \ce{CO}-Vergiftung
 &
\\\addlinespace
 &
Bläulich 
 &
Zyanose: Hypoxie!
 &
\\\addlinespace
\TabSub{Körperlich}
 &
Einziehungen an den Rippen
&
Einsatz der Atemhilfsmuskulatur
&
\\\addlinespace
&
Aufstützen
&
Einsatz der Atemhilfsmuskulatur
&
\\\addlinespace
&
Nasenflügeln \index{Nasenflügeln}
&
Atemnot, besonders bei Kleinkindern
&
\\\addlinespace
&
Aufrechte Position
&
&
\\\addlinespace
&
"`Paradoxe Atmung"'
&
Brustkorb senkt sich bei Einatmung
&
\\\addlinespace
\TabSub{Geräte}
&
Pulsoxymetrie
&
\FullPrettyRef{topic:pulsoxymetrie}
&
\\\addlinespace
&
\Z{Kapnometrie}
&
\Z{\ce{CO2}-Messung}
\\\addlinespace\bottomrule
\end{tabulary}
\end{WidePar}

\end{table}


\MassnahmenMK{Spezielle
Maßnahmen}{Insuffiziente
Atmung (AF \textless\,8 oder
\textgreater\,30\PerMin*,
bzw. AZV zu niedrig)}{m:atmung-insuffizient}
{\Ca{Vitale Bedrohung}
}{}

\MassnahmenMK{Spezielle
Maßnahmen}{Atemstillstand}{m:atemstillstand}{
\Cave{Achtung: Bei einem mit Atemstillstand vorgefundenen
Patienten ist grundsätzlich von einer Reanimation
auszugehen.}}{}
% \MassnahmenNG{Spezielle
% Maßnahmen}{Insuffiziente
% Atmung (AF \textless\,8 oder
% \textgreater\,30\PerMin*,
% bzw. AZV zu niedrig)}{\label{m:atmung-insuffizient}}
% {\Ca{Vitale Bedrohung}
% \begin{itemize}
%   \item \TxMassVit
%   \item Assistierte Beatmung
%   \item Ursachenforschung
% \end{itemize}
% }
% 
% \MassnahmenNG{Spezielle
% Maßnahmen}{Atemstillstand}{\label{m:atemstillstand}}{
% \Cave{Achtung: Bei einem mit Atemstillstand vorgefundenen
% Patienten ist grundsätzlich von einer Reanimation
% auszugehen.}
% 
% Vgl. die entsprechenden Kapitel im Teil Erste Hilfe und
% Sanitätshilfe bezüglich der genauen Durchführung der
% Maßnahmen.
% 
% \Z{Nur in sehr seltenen Fällen kann primär ein isolierter
% Atemstillstand ohne gleichzeitigem Kreislaufstillstand
% beobachtet werden. Auf diese Spezialfälle (z.\,B.
% Opiatintoxikation) wird hier nicht weiter eingegangen. Auch
% der iatrogene Atemstillstand (z.\,B. im Rahmen einer Narkose) wird an dieser Stelle nicht
% behandelt.}}








\clearpage
%%%%%%%%%%%%%%%%%%%%%%%%%%%%%%%%%%%%%%%%%%%%%%%%%%%%%%%%%%%
%%%%%%%%%%%%%%%%%%%%%%%%%%%%%%%%%%%%%%%%%%%%%%%%%%%%%%%%%%%
%%%%%%%%%%%%%%%%%%%%%%%%%%%%%%%%%%%%%%%%%%%%%%%%%%%%%%%%%%%
%%%%%%%%%%%%%%%%%%%%%%%%%%%%%%%%%%%%%%%%%%%%%%%%%%%%%%%%%%%
\section[Herz- / Kreislauf]{Herz-\,/\,Kreislauf}
\label{bkm:RefHeading42514202}

\subsection{Blut als Sauerstofftransporteur}

Durch das Kreislaufsystem wird sauerstoff- und
nährstoffreiches Blut gepumpt. Ausschlaggebend für die
Sauerstoffversorgung ist der \I{Sauerstoffgehalt} des
Blutes. Dieser wird wesentlich durch den Gasaustausch, aber
auch durch die Aufnahmekapazität des Blutes bestimmt.
Wieviel Sauerstoff das Blut aufnehmen kann, hängt vor allem
von der Anzahl der roten Blutkörperchen (\I{Erythrozyten})
bzw. der Konzentration des Blutfarbstoffes \I{Hämoglobin}
ab.

\Z{Im arteriellen Blut sind pro Liter ca 0,2\Liter*
Sauerstoff enthalten. Im Normalfall werden davon nur ca. 25\PC* vom
Körper aufgenommen, der Rest gelangt ungenutzt in das venöse
System zurück.
\cite{Silbernagl3-Kreislaufsystem}}

\subsection{Herzleistung}

Das Herz kontrahiert beim Erwachsenen mit einer
\I{Herzfrequenz} (\I{HF}) von \VonBis{60}{100}\PerMin* und wirft dabei jedes Mal ein
\I{Schlagvolumen} (\I{SV}) von ca. \VonBis{70}{80}\Ml* aus
\cite{Silbernagl3-Herz,Silbernagl3-Kreislaufsystem}. Daraus ergibt sich
\E{in Ruhe} ein \I{Herzminutenvolumen} (\I{HMV}) von ca.
5\,--\,6\,--\,8\,\nicefrac{\Liter}{\Min}.
Bei Belastung kann das Herzminutenvolumen beträchtlich
gesteigert werden.



\begin{minipage}{\linewidth}
\FormatTable
\Captionof{table}{Gegenüberstellung der
wichtigsten Kennzahlen, die die Herz- und Lungenleistung
beeinflussen.}
\begin{tabularx}{\linewidth}{XXX}\addlinespace\toprule\addlinespace
\noindent\TabHeading{~} & \noindent\TabHeading{Herz} &
\noindent\TabHeading{Lunge}
\\\addlinespace\midrule\addlinespace 
\TabSub{Was?} & Blut &
Luft \\\addlinespace
\TabSub{Qualität} & Sauerstoffgehalt\Z{\footnote{\Z{Nicht
gleichzusetzen mit der Sauerstoffsättigung}}} &
Sauerstoffkonzentration 
\\\addlinespace\midrule\addlinespace
\TabSub{Wie oft?}
& 
Herzfrequenz 

(HF) 
&
Atemfrequenz 

(AF)
\\\addlinespace
\TabSub{Wieviel?}

(pro Aktion)
&
Schlagvolumen

(SV)
&
Atemzugsvolumen

(AZV)
\\\addlinespace\midrule\addlinespace
\TabSub{Erzielte Menge}
&
Herzminutenvolumen 

(HMV)
&
Atemminutenvolumen

(AMV)
\\\addlinespace\bottomrule
\end{tabularx}
\end{minipage}

%%%%%%%%%%%%%%%%%%%%%%%%%%%%%%%%%%%%%%%%%%%%%%%%%%%%%%%%%%%
%%%%%%%%%%%%%%%%%%%%%%%%%%%%%%%%%%%%%%%%%%%%%%%%%%%%%%%%%%%
%%%%%%%%%%%%%%%%%%%%%%%%%%%%%%%%%%%%%%%%%%%%%%%%%%%%%%%%%%%
\subsection{Blutdruck} 
\label{topic:blutdruck}



% test:
% 
% \glsentrydesc{blutdruck}
% 
% :ende test

\Layout{0}{\TxBeschreibung}
{Durch die Pumpfunktion des Herzens entsteht ein
lebenswichtiger Druck im (arteriellen) Gefäßsystem, der
\I{Blutdruck}.
Da das Herz rhythmische kontrahiert, kommt es
infolge des Blutauswurfes zu synchronen -- ebenfalls
rhythmischen -- Druckveränderungen. Während des Auswurfs von
Blut in das Gefäßsystem (Herzanspannung, \I{Systole})
spricht man vom systolischen Blutdruck, während der Entspannungsphase des
Herzens (\I{Diastole}) vom diastolischen Blutdruck.
Natürlich ist der Blutdruck bei der Anspannung des Herzens höher als
bei der Entspannung (\E{Systolischer Blutdruck >
Diastolische Blutdruck)}.
}{
\begin{description}
    \item[systolisch:] Maximaldruck in der Arterie am Höhepunkt der
Kammerkontraktion
    \item[diastolisch:] Druck in der Arterie während die Kammer erschlafft ist
\end{description}
}{}{}{}


Zur Durchführung der Blutdruckmessung vgl.
\FullPrettyRef{topic:rr-messung}.




\begin{description}
    \item[Normbereich] zwischen \RRmmHg{90}{60} und
    \RRmmHg{140}{90} \index{Blutdruck!Normbereich}
    \item[Hypotonie] RR {\textless}~\RRmmHg{90}{60}
    \index{Hypotonie}
    \item[Hypertonie] RR {\textgreater}~\RRmmHg{140}{90}
    \index{Hypertonie}
\end{description}

{\tiny Werte beziehen sich auf den Erwachsenen.}

Achtung: Eine Hypertonie ist zwar eine Erkrankung, aber
nicht immer im Einsatz von Bedeutung. Viele Patienten sind
an ihren erhöhten Blutdruck gewohnt und haben keine
Beschwerden. Es darf nie nur nach dem Wert gehandelt werden,
es muss auch immer bedacht werden:

\begin{itemize}
    \item \EE{Zustand} des Patienten
    \item Der \EE{sonst übliche} Blutdruck des Patienten
    (erfragen, Blutdruckprotokoll)
    \item Der \EE{Verlauf} des Blutdruckes während der
    Behandlung (Traumapatienten!)
    \item Der \EE{Erregungszustand} des Patienten.
\end{itemize}

\minisec{Beispiele}

\emph{Ein 65-jähriger, männlicher Patient klagt über
Schwindelgefühl, sie messen einen Blutdruck von \RRmmHg{100}{70}. Im
Blutdruckprotokoll finden sie für die vergangenen vier Tage folgende Einträge:
\RR{170}{110}, \RR{140}{115}, \RR{175}{113}, \RR{183}{112}, \ldots}


Wie beurteilen Sie den Blutdruck?
\index{Blutdruck!Beurteilung}

\begin{quote}
Der Blutdruck ist für den Patienten zu niedrig, der Körper
ist auf eine so plötzliche Umstellung nicht vorbereitet
gewesen. Durch das Anamnesegespräch erfahren sie, dass er
heute das erste Mal einen "`Nitro-Spray"' verwendet hat. Sie
vermerken dies am Einsatzprotokoll. Im Krankenhaus erfahren
Sie, dass wahrscheinlich eine Überdosierung dieses
Medikaments die Ursache dieses Blutdruckabfalls war.
${\rightarrow}$ Obwohl der Patient laut Lehrbuch einen
"`schönen"' Blutdruck hatte, war dieser dennoch zu niedrig!
\end{quote}

\emph{Ein 56-jähriger Patient klagt über starke
kolikartige Schmerzen im linken Unterbauch. Sie messen einen Blutdruck von
\RRmmHg{155}{95}. Der Patient kennt seinen "`gewöhnlichen"' Blutdruck
nicht.}  

Wie beurteilen Sie den Blutdruck?

\begin{quote}
Der Blutdruck ist laut Lehrbuch erhöht, aber dies ist
erklärbar die starken Schmerzen und den damit verbundenen
Erregungszustand. Die Beschwerden stehen wahrscheinlich in
keinem ursächlichen Zusammenhang mit den Blutdruck.
\end{quote}

\emph{Eine 50-jährige Patientin hat Nasenbluten. Sie
messen einen Blutdruck von \RRmmHg{220}{120}. Auf Nachfrage gibt sie an
etwas Kopfsschmerzen zu haben, und etwas schwindlig sei sie auch.} 

Wie beurteilen Sie den Blutdruck? 

\begin{quote}
Der Blutdruck ist sehr stark erhöht, außerdem zeigt die
Patientin Zeichen einer hypertensiven Krise (Nasenbluten, Kopfweh, Schwindel).
Hier ist der erhöhte Blutdruck Ursache für die Beschwerden und muss
mittels Medikamente unbedingt gesenkt werden. 
\end{quote}


%%%%%%%%%%%%%%%%%%%%%%%%%%%%%%%%%%%%%%%%%%%%%%%%%%%%%%%%%%%
%%%%%%%%%%%%%%%%%%%%%%%%%%%%%%%%%%%%%%%%%%%%%%%%%%%%%%%%%%%
%%%%%%%%%%%%%%%%%%%%%%%%%%%%%%%%%%%%%%%%%%%%%%%%%%%%%%%%%%%
\subsection{Kreislaufstillstand}\label{bkm:RefHeading42514602}\label{topic:kreislaufstillstand}
\index{Kreislaufstillstand}

\Definition{Totales Pumpversagen des Herzens und somit
Stillstand des Kreislaufes.}

Das Pumpversagen kann entstehen durch:

\begin{itemize}
  \item Herzrhytmusstörungen\begin{itemize}
   \item Asystolie (keine Herztätigkeit)
    \item Kammerflimmern (unkoordinierte Herztätigkeit)
    \item Zu schnelle Herzfrequenz, bei der sich das Herz
    nicht mehr füllen oder etwas auswerfen kann ((Pulslose)
    Ventrikuläre Tachykardie) 
    \item Elektromechanische
    Entkopplung: keine
    Pumpreaktion auf Signale des Reizleitungssystems im Herz
\end{itemize}
 \item Mechanische Blockaden
 \item Blutverlust
\end{itemize}


\Layout{2}{Ursachen}
{
\begin{itemize}
  \item Minderversorgung des Herzmuskels mit \ce{O2}\par
  \begin{inparaitem}
    \item Herzinfarkt,
    \item Herzinsuffizienz, \ldots
  \end{inparaitem}
  \item Störungen der Atmung\par
  \begin{inparaitem}
    \item Asthma bronchiale,
    \item Lungenembolie,
    \item zentrale Atemregulationsstörungen, \ldots
  \end{inparaitem}
  \item Sonstige Störungen\par
  \begin{inparaitem}
    \item Trauma,
    \item Vergiftungen,
    \item thermische Schäden,
    \item Elektrolytverschiebungen, \ldots
  \end{inparaitem}
\end{itemize}
}{}{}{}{}

\Layout{0}{Folgen und Symptome}
{Unmittelbar anch Aussetzen des Herzschlages ist der Patient
\E{pulslos} und der Blutkreislauf steht still. Es wird im
schwindlig und schlecht, nach kurzer Zeit (ca. \VonBis{15}{20}\Sek*) wird er
\E{bewusstlos}. Nach weiteren 10\Sek* tritt ein
\E{Atemstillstand} ein (oder eine \I{Schnappatmung}). 

Nach \VonBis{3}{5}\Min* treten Hirnschäden auf, welche nicht
mehr reversibel sind. In kühler Umgebung kann diese
Zeitspanne etwas ausgedehnt werden.
Viele Patienten haben, auch nach erfolgreicher
Reanimation, im weiteren Verlauf neurologische Defizite. 
}{
\begin{itemize}
  \item Aussetzen des Herzschlages $\rightarrow$ 
  sofort pulslos!
  \item[\Sign{→}] Bewusstlosigkeit nach \VonBis{15}{20}\Sek*
  \item[\Sign{→}] Atemstillstand/Schnappatmung nach ca. 10\Sek*
  \item[\Sign{→}] \EE{Irreversible Hirnschäden nach \VonBis{3}{5}\Min*}
\end{itemize}

}{}{}{}




\Layout{0}
{\TxQuerverweise}
{\begin{itemize}
  \item Reanimation, Erste Hilfe:
  \FullPrettyRef{topic:reanimation-eh}
  \item Basic Life Support:
  \CO{[BLS]}, \FullPrettyRef{chp:BLS}
  \item Advanced Life Support: \CO{[ALS]},
  \FullPrettyRef{chp:ALS}
\end{itemize}}
{}
{}
{}
{}




%%%%%%%%%%%%%%%%%%%%%%%%%%%%%%%%%%%%%%%%%%%%%%%%%%%%%%%%%%%
%%%%%%%%%%%%%%%%%%%%%%%%%%%%%%%%%%%%%%%%%%%%%%%%%%%%%%%%%%%
%%%%%%%%%%%%%%%%%%%%%%%%%%%%%%%%%%%%%%%%%%%%%%%%%%%%%%%%%%%
\subsection{\Z{Herzinsuffizienz, Herzversagen}}
\index{Herzinsuffizienz}\index{Herzversagen}

Weitere Störungen des Herzens werden 
unter \FullPrettyRef{topic:herz-kreislauf-stoerungen} besprochen. 




%%%%%%%%%%%%%%%%%%%%%%%%%%%%%%%%%%%%%%%%%%%%%%%%%%%%%%%%%%%
%%%%%%%%%%%%%%%%%%%%%%%%%%%%%%%%%%%%%%%%%%%%%%%%%%%%%%%%%%%
%%%%%%%%%%%%%%%%%%%%%%%%%%%%%%%%%%%%%%%%%%%%%%%%%%%%%%%%%%%
%%%%%%%%%%%%%%%%%%%%%%%%%%%%%%%%%%%%%%%%%%%%%%%%%%%%%%%%%%%
\section{Schock}
\index{Schock}
\label{topic:schock}

\Definition{"`Lebensbedrohliche Kreislaufstörung aufgrund von
Unterversorgung des Körpers mit Blut und Sauerstoff infolge des
Versagens des Kreislaufsystems."'}


% \DefinitionExp{Schock}

\Fazit{Es können einzelne oder mehrere Teile des
Kreislaufsystems betroffen sein  (Blut, Gefäße, Herz).}

\Fazit{Ob ein Schock oder eine Schockgefahr besteht wird anhand von
\underline{Schocksymptomen} und anhand der
\underline{Verdachtsdiagnose} beurteilt.}

% TODO Graphik Schockentwicklung RR und HF



\noindent Schock bedeutet:

\begin{itemize}
  \item Akute Kreislaufinsuffizienz =
  \EE{Missverhältnis
  Herzauswurf\,:\,Sauerstoffbedarf} des Körpers
  \item Mögliche Ursachen:
  \begin{enumerate}
    \item Herz pumpt zu wenig
    \item Zu wenig Blut bzw. Flüssigkeit vorhanden
    \item Gefäße sind zu weit gestellt und Blut versackt
  \end{enumerate}
\end{itemize}

\Fazit{Im Endeffekt kommt immer zu wenig Blut mit
Sauerstoff zu den Organen. Dadurch ist die
Gewebedurchblutung für die Versorgung aller Organe mit
Sauerstoff nicht ausreichend.}

Körpereigene Mechanismen können versuchen gegenzusteuern um
den Schock zu kompensieren. Dies geht allerdings nur eine
gewisse Zeit und in einem gewissen Ausmaß gut, daher ist
eine schnell  einsetzende Therapie wichtig! Bei länger
bestehendem Schock kann dieser Zustand irreversibel werden.

\E{Kein} Schock besteht bei kurzzeitigen 
Bewusstseinsstörungen wie Kollaps/Synkope
(\underline{kurzfristige} Blutverteilungsstörung,
kurzfristige Minderdurchblutung des Gehirns,  flüchtige Symptome)

\Layout{0}{Kompensation}
{
\index{Schock!Kompensation}
Der Körper versucht selber dagegen etwas zu tun und gegenzusteuern
("`kompensieren"'). Der
wichtigste Mechanismus ist die
\index{Zentralisation}\T{Zentralisation}. Darunter versteht
man die Blutumverteilung zugunsten lebenswichtiger Organe wie
Herz, Gehirn, Lunge und Leber. Dabei werden z.\,B. Haut, Muskel
und Darm
sowie die Extremitäten (=Peripherie) kaum mehr  durchblutet.
Periphere Venen kollabieren und die Rekapillarisierungszeit
verlängert sich. Die Hautdurchblutung wird verringert und
die Haut wird kühl und bleich.

Wenn die Kompensationsmaßnahmen 
nicht mehr ausreichen, spricht man vom \E{"`dekompensierten
Schock"'}\index{Schock!dekompensierter}. Es besteht
spätestens dann akute \EE{Lebensgefahr}!
}{
\begin{itemize}
    \item "`Gegensteuerung"'
    \item Zentralisation
    \item Ende: "`Dekompensierter Schock"'
\end{itemize}
}{}{}{}




% \Cave{Bei jedem Patienten ist zu erheben, ob ein Schock
% oder eine realistische Schockgefahr besteht oder nicht bzw. festzustellen um welche Art von
% Schock es sich handelt (s.\,u.)!}

\Cave{Ein Schockzustand ist grundsätzlich als vitale
Bedrohung einzustufen!}

\Layout{2}{Schock beim Kind}
{Die Schockschere gilt beim Kind nicht, da die
Normalwerte andere sind! Kinder bleiben außerdem lange
kreislaufstabilstabil, verfallen dann aber schlagartig!

}{}
{}
{}
{}
{}

%%%%%%%%%%%%%%%%%%%%%%%%%%%%%%%%%%%%%%%%%%%%%%%%%%%%%%%%%%%
%%%%%%%%%%%%%%%%%%%%%%%%%%%%%%%%%%%%%%%%%%%%%%%%%%%%%%%%%%%
%%%%%%%%%%%%%%%%%%%%%%%%%%%%%%%%%%%%%%%%%%%%%%%%%%%%%%%%%%%
\subsection{Allgemeine Schocksymptomatik}
\label{topic:schock-allgemeine-symptomatik}
\index{Schock!allgemeine Symptomatik}

Die Symptome eines Schocks sind abhängig vom jeweiligen
Stadium. Problematisch ist, dass am Anfang die Symptome
eher diskret sind. 
\Fazit{Das erste Anzeichen eines beginnenden Schocks ist
oft eine kühle
Haut\cite{Kaplan2001}!}

Natürlich kann es
auch andere Ursachen für eine kühle Haut geben, z.\,B. eine
kalte Umgebungstemperatur.

Die Tabelle
\prettyref{tab:schocksymptomatik} gibt eine entsprechende Übersicht.

\begin{table}[H]
\begin{WidePar}
\Captionof{table}{Allgemeine
Schocksymptomatik.\label{tab:schocksymptomatik}}
\CorrectionItemizeBox\FontTableBase

\begin{tabularx}{\linewidth}{p{3.5cm}m{1,2cm}X}
\addlinespace\toprule\addlinespace
\noindent\TabHeading{Stadium}
&
\noindent\TabHeading{Pat.}
&
\noindent\TabHeading{Symptome}
\\\addlinespace\midrule\addlinespace
{\sffamily\Large\color{DarkGreen} \index{Schockgefahr}Schockgefahr}
&

\centering\includegraphics[width=1cm]{./icons/sm-basic.png}
&
 \begin{itemize}
     \item Keine eindeutigen Anzeichen
     \item Vorausschauend handeln!
     \item Verdachtsdiagnose bedenken.  Muss man damit
     rechnen, dass sich ein Schock entwickelt?  Besteht
     Schockgefahr?
 \end{itemize}
\\\addlinespace
{\sffamily\Large\color{OrangeRed}Schock}
&

\centering\includegraphics[width=1cm]{./icons/sm-pale.png}
&
 \begin{itemize}
     \item HF ${\uparrow}$
     \item Mäßiger RR-Abfall, aber RR noch {\textgreater} HF
     \item Neurologisch unauffällig, ev. euphorisch oder
     agitiert
     \item Haut blass, evtl. kaltschweißig
     \item Venen sind kollabiert oder gestaut
     \item Peripherie schlecht durchblutet, Rekapzeit erhöht
 \end{itemize}
\\\addlinespace
{\sffamily\Large\color{red}Schwerer Schock

Vollbild}
&

\centering\includegraphics[width=1cm]{./icons/sm-pale-frown.png}
&
\begin{itemize}
  \item Tachykardie (HF ${\uparrow}$, "`fliehender"' Puls)
  \item RR-Abfall (schwacher Puls)
  \begin{itemize}
    \item Herzfrequenz {\textgreater} systolischer RR
  \end{itemize}
  \item Tachypnoe (schnelle, flache Atmung)
  \item Neurologisch auffällig, getrübt
  \item Haut blass, zyanotisch (bläulich), feucht
  \item Venen kollabiert oder gestaut
  \item Nur mehr schwer behebbar! ${\rightarrow}$ Schock  möglichst
  frühzeitig erkennen \&  behandeln
\end{itemize}
\\\addlinespace
{\sffamily\Large\color{LightSlateGray}Endstadium}
&

\huge\centering $\skull$
&
\begin{itemize}
  \item Haut grau, klebrig
  \item Versagen der Herzfunktion
  \item RR kaum mehr messbar
  \item Evtl. Koma
  \item Evtl. Schnappatmung
  \item Patient ist dabei zu sterben.
\end{itemize}
\\\addlinespace\bottomrule
\end{tabularx}
\end{WidePar}
\end{table}


\clearpage % TEMP temp clearpage
\Layout{0}{Schockschere}
{
\index{Schockschere}
\label{topic:schockschere}
Normalerweise ist die Herzfrequenz niedriger als der
systolische Blutdruck. (Achtung: \EE{Kinder haben andere Normalwerte für HF und
RR!} Die Faustregel ist deshalb nicht anwendbar!) Wenn der Körper versucht den Schock
zu kompensieren, kann sich dieses Verhältnis umkehren. Man
spricht dann von der Schockschere, der Patient befindet sich
dann \EE{bereits  im Vollbild eines Schocks!} 

\E{Die
Schockschere ist ein \EE{spätes Zeichen!} Auch wenn sich ein
Patient (noch) nicht in der Schockschere befindet, kann er
trotzdem einen Schock haben!}



% % TODO Schockschere
% \Captionof{figure}{Schockschere. \AuthNotes{Achtung! Bild
% fehlt!}}

\Fazit{Wenn sich das Verhältnis 'Herzfrequenz zu Systolischem
Blutdruck' umkehrt, befindet sich der Patient bereits in
einem schweren Schock.} 
\Fazit{Die Schockschere ist ein \EE{Spätzeichen!} Auch wenn
die Herzfrequenz niedriger ist als der systolische Blutdruck
kann sich der Patient bereits im Schock befinden!}
\Fazit{Ob ein Schock oder eine Schockgefahr besteht, wird anhand von
\E{Schocksymptomen} und anhand der
\E{Verdachtsdiagnose} beurteilt, \EE{\E{nicht}} primär anhand der Schockschere!}
}{
\begin{itemize}
  \item RR\textsubscript{syst}\,{\textless}\,HF
  \item Dann bereits \E{schwerer} Schock
  \item Spätzeichen, auch ohne Schockschere kann Schock bestehen.
  \item[\Ca{!}] \EE{Schockbeurteilung}: Schocksymptome und
  Verdachtsdiagnose, \E{nicht} primär Schockschere!
  \item Für Kinder nicht anwendbar
\end{itemize}
}{}{}{}




%%%%%%%%%%%%%%%%%%%%%%%%%%%%%%%%%%%%%%%%%%%%%%%%%%%%%%%%%%%
%%%%%%%%%%%%%%%%%%%%%%%%%%%%%%%%%%%%%%%%%%%%%%%%%%%%%%%%%%%
%%%%%%%%%%%%%%%%%%%%%%%%%%%%%%%%%%%%%%%%%%%%%%%%%%%%%%%%%%%
\subsection{Schock -- Arten}
\index{Schock!Arten}
\label{topic:schock-arten}

Je nach dem zugrunde liegenden Mechanismus unterscheidet man
verschiedene Schockarten:

\begin{labeling}[~:]{Volumenmm}
  \item[Volumen] Es ist \E{zu wenig Flüssigkeit} bzw. Blut
  im Gefäßsystem vorhanden
  \begin{itemize}
    \item \EE{Hypovolämischer} Schock
  \end{itemize}
  \item[Herz] Der \E{Herzauswurf} (bei normalem Volumen) ist
  zu gering.
  \begin{itemize}
    \item \EE{Kardiogener} Schock
  \end{itemize}
  \item[Gefäße] Die Gefäße sind inadäquat \E{weitgestellt},
  das Blut versackt
  \begin{itemize}
    \item \EE{Anaphylaktischer} Schock
    \item \EE{Septischer} bzw. \EE{toxischer} Schock
    \item \EE{Neurogener} Schock
  \end{itemize}
\end{labeling}



\subsubsection{Hypovolämischer Schock}
\label{topic:schock-hypovolaemischer}
\index{Schock!Hypovolämischer}

\Definition{Schock aufgrund eines Verlustes von Blutvolumen aus den
Gefäßen.}




Es kann dabei verloren gehen: 

\begin{itemize}
    \item Blut
    \item Plasma
    \item Wasser und Elektrolyte
\end{itemize}

\paragraph{Ursachen}




\begin{center}
\CorrectionItemizeBox\FontTableBase
\begin{tabularx}{\linewidth}[H]{XX}
\addlinespace\toprule\addlinespace
Flüssigkeitsverlust nach außen
&
Flüssigkeitsverlust nach innen
\\\addlinespace\midrule\addlinespace
\begin{itemize}
  \item Blutung
  \item Gastrointestinal (Durchfall, Diarrhoe)
  \item Verbrennungen (Hautbarriere versagt  groß\-flächig)
  \item Renal (Diabetes mellitus, Diuretika, Polyurie bei 
  Nierenversagen)
\end{itemize}
&
\begin{itemize}
  \item Innere Blutungen
  \begin{itemize}
    \item Nach Trauma, rupturierte Eileiter\-schwanger\-schaft, etc
  \end{itemize}
  \item Wassereinlagerung ins Gewebe
  \item \Z{Aszites
  ("`Wasserbauch"' bei
  Lebererkrankungen  etc.)}
\end{itemize}
\\\addlinespace\bottomrule
\end{tabularx}
\end{center}
Bei akutem Flüssigkeitsverlust kann ein Verlust von bis zu
20\PC* des Blutvolumens  gut kompensiert werden. Auf größere
Verluste  folgt ein Absinken von Herzauswurf und RR!

% \MassnahmenNG{Spezielle Maßnahmen}{Hypovolämischer Schock}{}
% {%
% \begin{itemize}
%   \item Allgemein: \EE{Allgemeine Maßnahmen der
%   Schockbekämpfung} (s.
%   \prettyref{m:am-schock})
%   \item Blutstillung  (weitere Verluste verhindern!)
%   \item Ggfs. Trauma-Vorgehen beachten
%   \item Wenn möglich schonende Bergung
%   \item Lagerung: grundsätzlich Beine hoch
%   
%   Wichtige Ausnahmen: {\textperiodcentered}\,kardiale Ursache,
%   {\textperiodcentered}\,Kopf- / Brust\-ver\-letzungen,
%   {\textperiodcentered}\,Hirndruck (SHT),
%   {\textperiodcentered}\,Knochenbrüche der Beine oder des
%   Beckens,
%   {\textperiodcentered}\,Wirbelsäulenverletzungen,
%   {\textperiodcentered}\,Bewusstlosigkeit ${\rightarrow}$ sonst Flachlagerung oder stabile Seitenlage \A{KOCH: darüber müssen wir reden: Schock vor Schädel?}
%   %TODO Schock vor Schädel
%   \item Evtl. Aviso (Vorankündigung) im  Krankenhaus /
%   Schockraum
%   \item Abtransport: rasch, schonend, nicht überstürzt
% \end{itemize}
% }

\AuthNotes{Begriff "`Schocklagerung"'
ein für alle mal in der Ausbildung streichen, nur Hinweis
dass dieser Ausdruck veraltet ist sollten sie ihn wo
aufschnappen. Stattdessen "`Beine-hoch"'-Lagerung.

Wenn KI gegen Beine hoch: Flachlagerung bzw. stab. Seitenlage empfehlen?

STEU: ja, Flachlagerung wenn KI. 2009-06-28}


\subsubsection{Kardiogener Schock}
\index{Schock!Kardiogener}
\label{topic:schock-kardiogener}

\Definition{Schock infolge mangelndem Herzauswurfes durch
insuffiziente Pumpleistung.} 

\Layout{0}{\TxBeschreibung}{%
Zu einem kardiogenen Schock kann es bei einem \EE{akut
auftretenden Herzversagen}, z.\,B. durch einem akuten
\E{Myokardinfarkt} (häufigste Ursache) oder
\E{Herzrhythmusstörungen}, kommen. Die \E{Lungenembolie} als
Auslöser eines Schocks wird auch zum kardiogenen Schock
gerechnet, da es durch die Verstopfung des kleinen
Kreislaufes zu einer mangelhaften Füllung des Herzens und
damit zu einer mangelhaften Auswurfleistung kommt.
}{%
\begin{itemize}
  \item Herzversagen, Herzinsuffizienz
  \begin{itemize}
    \item Myokardinfarkt (häufigste Ursache)
    \item Herzrhythmusstörungen
  \end{itemize}
  \item Lungenembolie
\end{itemize}
}{}{}{}

\Layout{0}{\TxSymptome}{%
Typisch ist ein Blutdruckabfall. Durch den
\EE{Blutrückstau} kann es zu Stauungszeichen kommen
(Lungenödem bei Linksherzversagen, Beinödeme und gestaute
Halsvenen bei Rechtsherzversagen). Häufig klagen die
Patienten über \EE{Atemnot}.

\vspace{2\baselineskip}

}{%
\begin{itemize}
  \item RR-Abfall
  \item Stauungszeichen
  \begin{itemize}
    \item Lungenödem
    \item Beinödeme, gestaute Halsvenen
  \end{itemize}
  \item Atemnot
\end{itemize}
}{}{}{}

% \begin{center}
% \CorrectionItemizeBox\FontTableBase
% \begin{tabularx}{\linewidth}[H]{XX}
% \toprule
% Ursachen &
% charakterisiert durch
% \\\midrule
% \begin{itemize}
%   \item Herzversagen, Herzinsuffizienz
%   \begin{itemize}
%     \item Myokardinfarkt (häufigste Ursache)
%     \item Herzrhythmusstörungen
%   \end{itemize}
%   \item Lungenembolie
% \end{itemize}
% &
% \begin{itemize}
%   \item RR-Abfall
%   \item Eingeschränkte Pumpleistung
%   \begin{itemize}
%     \item Blut staut sich zurück
%   \end{itemize}
% \end{itemize}
% \\\bottomrule
% \end{tabularx}
% \end{center}



% \MassnahmenNG{Spezielle Maßnahmen}{Kardiogener Schock}{}{
% 
% \begin{itemize}
%   \item Allgemein: \EE{Allgemeine Maßnahmen der
%   Schockbekämpfung}
%   (\prettyref{m:am-schock})
%   \item \TxMassVitBes
%   \begin{itemize}
%     \item \ce{O2}: 10\LPerMin*{}   oder mehr
%     \item Lagerung: Oberkörper hoch! Evtl. Beine hängen  lassen.
%   \end{itemize}
%   \item Medikamentöse Therapie durch Notarzt
%   \Z{(Diuretika, Katecholamine, Heparin)}
% \end{itemize}
% }







\subsubsection{Anaphylaktischer Schock}
\label{topic:schock-anaphylaktischer}
\index{Schock!Anaphylaktischer}
\index{Schock!allergischer}

\Definition{Schock aufgrund einer überschießenden
allergischen Reaktion -- "`Allergischer
Schock"'.}



\Layout{0}{\TxBeschreibung}{%
Ein anaphylaktischer Schock wird durch eine
über\-schieß\-ende Reaktion der Körperabwehr (Immunreaktion)
auf bestimmte Stoffe (Antigene) ausgelöst.

Durch von den Immunzellen gebildete Stoffe werden
die Gefäße erweitert, zusätzlich steigt die
Durchlässigkeit der Gefäße\ZFootnote{(Gefäßpermeabilität)}.
Das Blut versackt und es kommt zu einem  Plasmaverlust in
den Zwischengewebsraum\ZFootnote{(Interstitium)}; der
Blutdruck fällt massiv ab!
}{%
\begin{itemize}
  \item Überschießende
  Reaktion der Körperabwehr
  \item Gefäße erweitert, Blut versackt
  \item Durchlässigkeit der Gefäße ↑, Plasmaverlust in
  Zwischengewebsraum
\end{itemize}
}{}{}{}


\Layout{0}{Mögliche Auslöser}
{%
\begin{itemize}
    \item Pollen
    \item Medikamente, Antibiotika
    \item Kontrastmittel (für radiologische Untersuchungen)
    \item tierische Gifte (Bienenstich etc.)
    \item Lebensmittel (Meeresfrücjte, Nüsse, etc.)
    \item und vieles andere \ldots
\end{itemize}
}{%
\begin{itemize}
    \item Pollen
    \item Medikamente
    \item Kontrastmittel 
    \item tierische Gifte 
    \item u.\,v.\,a.\,m. 
\end{itemize}
}{}{}{}


\Layout{0}{\TxSymptome}
{
\prettyref{tab:anaphylaxie-stadien}.
}
{
\prettyref{tab:anaphylaxie-stadien}.
}
{}{}{}
% TODO Flush erklären


\begin{table}[H]
\CorrectionItemizeBox\FontTableBase
\Captionof{table}{Stadien der
Anaphylaxie\label{tab:anaphylaxie-stadien}}
\begin{tabulary}{\linewidth}[H]{cL}\addlinespace\toprule\addlinespace
\noindent\TabHeading{\mbox{Stadium}}
 &
\noindent\TabHeading{{Symptome}}
\\\addlinespace\midrule\addlinespace
\TabSub{I}
 &
Hautreaktion: evtl. Ausschlag, Rötung;
"`Quaddeln"', Flush, Ödeme

Allgemeinsymptome: Juckreiz, Zittern, Schwindel, Kopfschmerz

\\\addlinespace
\TabSub{II}

 &
RR ${\downarrow}$, Tachykardie

Übelkeit, Erbrechen, Durchfall 

\\\addlinespace
\TabSub{III}
 &
Allgemeine Schocksymptome

Bronchospasmus
\\\addlinespace
\TabSub{IV}
 &
Atem-/Kreislaufstillstand

Schwerer Schock
\\\addlinespace\bottomrule
\end{tabulary}
\end{table}

% \MassnahmenNG{Spezielle Maßnahmen}{Anaphylaktischer
% Schock}{}{%
% \begin{itemize}
%   \item Allgemein: \EE{Allgemeine Maßnahmen der
%   Schockbekämpfung}
%   (\prettyref{m:am-schock})
%   \item Antigen  nach Möglichkeit entfernen!
%   \item Lagerung: situationsgerecht.
%   Grundsätzlich Beine hoch, bei
%   Atemproblemen jedoch stark erhöhten
%   Oberkörper erwägen.
% \end{itemize}
% }


\subsubsection{Septischer Schock und Toxischer Schock}
\index{Schock!septischer}
\index{Schock!toxischer}
\label{topic:schock-septischer}
\label{topic:schock-toxischer}
Siehe auch \FullPrettyRef{topic:sepsis}.

\Definition{Schock infolge einer schweren, körperweiten Infektion
oder aufgrund eines, die Gefäße beeinflussenden, Giftstoffes
(Toxins).}


\Layout{0}{Ursachen}
{%
\index{Toxine!Toxischer Schock}%
Ein septisch-/toxischer
Schock wird durch diverse Bakterien und Toxine verursacht.
Die Erreger können auf unterschiedliche Wege in die Blutbahn
gelangen. Oft liegt ein chronischer Entzündungsherd vor, von
dem aus Keime in die Blutbahn streuen.  Mit Keimen
besiedelte und im Körper gelegene Fremdkörper, darunter
fallen auch medizinische Materialien wie
lang liegende Venenverweilkanülen, Harnkatheter u.\,ä., aber
auch Hygieneartikel wie vergessene Tampons, sind ideale
Brutstätten, da hier das Immunsystem des Körpers keinen Zugang hat. 

Durch die bakterielle oder toxische Gefäßwandschädigung kommt es
zum Flüssigkeitsaustritt in den Zwischengewebsraum sowie zu Gefäßweitstellung.
}{
\begin{itemize}
  \item Langanhaltender Entzündungsherd der plötzlich Erreger
  "`ins Blut streut"' und "`explodiert"'
  \item Fremdkörper (Venenverweilkanülen, \ldots)
  \item Tampons
  \item \ldots
\end{itemize}
}{}{}{}



\Layout{0}{\TxSymptome}
{%
Man kann die allgemeinen Schocksymptome wahrnehmen. Der
Patient kann sowohl hohes  oder auch kein Fieber haben
Die Haut ist anfänglich heiß, trocken und rötlich, später
ist sie eher bleich, grau oder sogar marmoriert.
}{
\begin{itemize}
  \item Schock
  \item Hohes oder kein Fieber
  \item Haut heiß, trocken, rot, später bleich, grau,
  marmoriert
\end{itemize}
}{}{}{}








% \MassnahmenNG{Spezielle Maßnahmen}{Septischer und Toxischer Schock}{}{
% \begin{itemize}
%   \item Allgemein: \EE{Allgemeine Maßnahmen der
%   Schockbekämpfung} (\prettyref{m:am-schock})
%   \item Lagerung: situationsgerecht, Beine hoch
%   \item Schutz vor Unterkühlung/Überwärmung
%   \item Aviso im Krankenhaus / Intensivstation
%   ("`Überwachungsbett"')
%   \item benötigt Antibiotika (im KH)
% \end{itemize}
% }

\subsubsection{Neurogener Schock}
\index{Schock!Neurogener}
\index{Schock!Spinaler}

\Definition{Schockzustand ausgelöst durch eine
Unterbrechung oder Störung der nervlichen Versorgung der
Kreislaufregulation. Es kommt zu inadäquater arterieller
und venöser Gefäßerweiterung.
\Z{(Vasodilatation)}}

\Layout{0}{Ursachen}
{Auch die Muskulatur der Blutgefäßwand wird mit Nerven
versogrt, damit die Weite der Gefäße reguliert werden kann.
Im Rahmen von Erkrankungen und Verletzungen kann diese
nervale Versorgung unterbochen sein, in weiterer Folge
weiten sich die betroffenen Gefäße und das Blut versackt. 
Typischerweise kommt dies im Rahmen einer
\EE{Wirbelsäulenverletzung} vor, bei der das Rückenmark
geschädigt bzw. unterbrochen wird (\I{Querschnitt-Syndrom}).
}{
\begin{itemize}
    \item Rückenmarksverletzung (Querschnitt-Syndrom)
    \item  Gefäßwandmuskulatur bekommen keine Nerven-Impulse
    und kontrahiert sich nicht mehr
    \item Gefäße weiten sich ${\rightarrow}$ Blut versackt
\end{itemize}
}{}{}{}





% \MassnahmenNG{Spezielle Maßnahmen}{Neurogener Schock}{}{
% \begin{itemize}
%   \item Allgemein: \EE{Allgemeine Maßnahmen der
%   Schockbekämpfung} (\prettyref{m:am-schock})
%   \item Lagerung: Flachlagerung (WS-Verletzung ist Kontraindikation gg.
%   Beine hoch!)
% \end{itemize}
% }

\clearpage % TEMP temp clearpage
\subsection{Maßnahmen der Schockbekämpfung}
\label{topic:schockbekaempfung}

\Fazit{Grundsatz: \EE{Der Schock muss bekämpft werden
bevor man ihn bemerkt.}}

\MassnahmenMK{Spezielle Maßnahmen}{Schock}
{m:am-schock}
{
}{}
% \MassnahmenNG{Allgemeine Maßnahmen}{Schockbekämpfung}
% {\label{m:am-schock}}
% {
% \begin{enumerate}
%   \item \TxMassVitBes
%   \begin{itemize}
%     \item Situationsgerechte  Lagerung, je nach Diagnose und Schockart
%     \item Sauerstoffgabe hochdosiert (ab 8\LPerMin*{} bei Schockgefahr,
%     15\LPerMin*{} oder mehr bei bestehendem Schock, je nach Schwere)
%     \item Monitoring: Besonders auf Haut, RR und HF achten!
%   \end{itemize}
%   \item Beengende Kleidung öffnen
%   \item Wärmeerhalt besonders beachten!
%   \item Spezielle Maßnahmen je nach Schockart bzw. Diagnose
% \end{enumerate}
% }
%     \AuthNotes{Freigabe: Sauerstoffgabe hochdosiert (ab 6l/min bei
%     Schockgefahr, 10l/min oder mehr bei bestehendem Schock, je nach Schwere}

% \MassnahmenWide{Spezielle Maßnahmen}{Schock}
% {\label{m:am-schock}}
% {\minisec{Allgemein}
% \begin{enumerate}
%   \item \TxMassVitBes
%   \begin{itemize}
%     \item Situationsgerechte  Lagerung, je nach Diagnose und Schockart
%     \item Sauerstoffgabe hochdosiert (ab 8\LPerMin*{} bei Schockgefahr,
%     15\LPerMin*{} oder mehr bei bestehendem Schock, je nach Schwere)
%     \item Monitoring: Besonders auf Haut, RR und HF achten!
%   \end{itemize}
%   \item Beengende Kleidung öffnen
%   \item Wärmeerhalt besonders beachten!
%   \item Spezielle Maßnahmen je nach Schockart bzw. Diagnose
%   \item Transportentscheidung: Vorankündigung
%   (Aviso); Spezialbett, je nach Schockart/Ursache
% \end{enumerate}
% 
% \minisec{Speziell}\CorrectionItemizeBox
% 
% \begin{tabularx}{\linewidth}{XXXXX}\toprule\addlinespace
% \multicolumn{5}{c}{\noindent\TabHeading{Schockart}}
% \\\addlinespace
% \noindent\TabHeading{Volumenmangel}
% &
% \noindent\TabHeading{Kardiogener}
% &
% \noindent\TabHeading{Anaphylaktischer}
% &
% \noindent\TabHeading{Septischer, Toxischer}
% &
% \noindent\TabHeading{Neurogener}
% \\\addlinespace\midrule\addlinespace
% \multicolumn{5}{c}{\TabSub{Lagerung}}
% \\\addlinespace\midrule\addlinespace
% Grundsätzlich Beine hoch
% 
% Kontraindikationen beachten!\footnote{Wichtige
% Ausnahmen: {\textperiodcentered}\,kardiale Ursache, {\textperiodcentered}\,Kopf- / Brust\-ver\-letzungen,
% {\textperiodcentered}\,Hirndruck (SHT),
% {\textperiodcentered}\,Knochenbrüche der Beine oder des
% Beckens,
% {\textperiodcentered}\,Wirbelsäulenverletzungen,
% {\textperiodcentered}\,Bewusstlosigkeit ${\rightarrow}$
% sonst Flachlagerung oder stabile Seitenlage \A{KOCH: darüber
% müssen wir reden: Schock vor Schädel?}}
% %TODO Schock vor Schädel
% &
% Oberkörper hoch! Evtl. Beine hängen  lassen.
% &
% Situationsgerecht: Grundsätzlich Beine
% hoch, bei Atemproblemen jedoch stark erhöhten
% Oberkörper erwägen.
% &
% situationsgerecht, Beine hoch
% &
% Flachlagerung 
% 
% (Cave WS-Verletzung!)
% \\\addlinespace\midrule\addlinespace
% \multicolumn{5}{c}{\TabSub{Sonstige Maßnahmen}}
% \\\addlinespace\midrule\addlinespace
% Blutstillung  (weitere Verluste verhindern!)
% 
% Wenn möglich schonende/schnelle Bergung\A{Schonend
%   oder schnell?}
% &
% Medikamentöse Therapie durch Notarzt\ZFootnote{(Diuretika, Katecholamine, Heparin)}
% &
% Auslösendes Antigen nach Möglichkeit entfernen!
% &
% 
% &
% \\\addlinespace\midrule\addlinespace
% \multicolumn{5}{c}{\TabSub{Transportentscheidung}}
% \\\addlinespace\midrule\addlinespace
% Schockraum, 
% 
% zügiger Abtransport
% &
% Ursachenabhängig
% &
% Intensivstation
% &
% Intensivstation
% \\\addlinespace\bottomrule
% \end{tabularx}
% }



\vspace{2em}
\clearpage % TEMP temp clearpage
%%%%%%%%%%%%%%%%%%%%%%%%%%%%%%%%%%%%%%%%%%%%%%%%%%%%%%%%%%%
%%%%%%%%%%%%%%%%%%%%%%%%%%%%%%%%%%%%%%%%%%%%%%%%%%%%%%%%%%%
%%%%%%%%%%%%%%%%%%%%%%%%%%%%%%%%%%%%%%%%%%%%%%%%%%%%%%%%%%%
%%%%%%%%%%%%%%%%%%%%%%%%%%%%%%%%%%%%%%%%%%%%%%%%%%%%%%%%%%%
\section{Temperaturregulation}
\label{topic:temperaturregulation}
\index{Temperatur!-regulation}

Die normale Körper-Kerntemperatur beträgt
\EE{37\textcelsius}. Verschiedene Mechanismen sorgen für
Wärmeerzeugung und -abgabe:


\begin{table}[H]\FormatTable
\begin{tabularx}{\linewidth}{XX}\addlinespace\toprule\addlinespace
\noindent\TabHeading{Wärmeproduktion durch}
&
\noindent\TabHeading{Wärmeabgabe über}
\\\addlinespace\midrule\addlinespace
\begin{itemize}
  \item Zellaktivität
  \item Stoffwechselvorgänge
  \item Bewegung (Muskelatmung)
\end{itemize}
&
\begin{itemize}
  \item Haut
  \item Schweiß
  
  \begin{itemize}
    \item Wirkung durch Verdunstungskälte
    \item Ca. 0,5\,{\FontUnits\nicefrac{\Liter}{d}}
    \item Bis max. 1,5\,{\FontUnits\nicefrac{\Liter}{h}} bei
    schwerer körperlicher  Arbeit
  \end{itemize}
\end{itemize}
\\\addlinespace\bottomrule
\end{tabularx}
\end{table}




%%%%%%%%%%%%%%%%%%%%%%%%%%%%%%%%%%%%%%%%%%%%%%%%%%%%%%%%%%%
%%%%%%%%%%%%%%%%%%%%%%%%%%%%%%%%%%%%%%%%%%%%%%%%%%%%%%%%%%%
%%%%%%%%%%%%%%%%%%%%%%%%%%%%%%%%%%%%%%%%%%%%%%%%%%%%%%%%%%%
%%%%%%%%%%%%%%%%%%%%%%%%%%%%%%%%%%%%%%%%%%%%%%%%%%%%%%%%%%%
\section{Stoffwechsel}
\label{topic:stoffwechsel}
\index{Stoffwechsel}
\index{Zuckerstoffwechsel}\index{Kohlenhydrate}\index{Fettstoffwechsel}
\index{Eiweißstoffwechsel}\index{Proteine}\index{Enzyme}

\Definition{Abläufe von Stoffumsetzung und 
Energiegewinnung, Nährstoffe werden in Wärme und Energie  umgewandelt}


\Layout{0}{\TxBeschreibung}
{% 
Der menschliche Körper hat drei wesentliche
Stoffwechselsysteme, die ihm zur Verfügung stehen. Man
unterscheidet den \EE{Zucker- bzw. Kohlehydratstoffwechsel},
den \EE{Fettstoffwechsel} und den \EE{Eiweißstoffwechsel}:

\begin{itemize}
  \item \EE{Zuckerstoffwechsel}: Zucker ist ein potenter
  \E{Energielieferant} und kann vom Körper schnell umgesetzt werden. Er wird durch
  Aufspaltung von komplexen \E{Kohlehydraten} (Stärke)
  gewonnen und zirkuliert in Form des \EE{Blutzuckers} in der
  Blutbahn. Die Leber speichert eine Zuckerreserve von ca.
  24\Hour*.
  
  Das Hormon \T{Insulin}\label{topic:insulin}, welches in der
  Bauchspeicheldrüse (Pankreas) produziert wird,  ermöglicht
  die \E{Aufnahme von Blutzucker in die Zellen} und \E{senkt
  dabei den Blutzuckerspiegel}.
  \item \EE{Fettstoffwechsel}: Fette haben einen doppelt so
  hohen Energiegehalt als Kohlenhydrate, können aber nicht so
  rasch vom Körper verwendet werden. Sie dienen außerdem
  u.\,a. als Lösungsmittel für Vitamine, Bausteine der
  Zellmembran und von Hormonen. \Z{(Weiters unterscheidet man
  zwischen gesättigten und ungesättigten Fettsäuren.)}
  \item \EE{Eiweißstoffwechsel}: Eiweiße (\T{Proteine})
  sind vergleichbar mit LEGO{\texttrademark}-Bausteinen: Sie
  sind aus Bausteinen zusammengebaut, den \Z{Aminosäuren}.
  Zusammengebaute Proteine können wieder beim
  Verdauungsvorgang \E{zerlegt} werden, mit den
  dadurch entstehenden Aminosäuren können wiederum andere,
  \E{neue Proteine} gebaut werden. Proteine sind in vielen Bereichen im menschlichen Körper
  unabkömmlich, sie kommen zum Beispiel als \I{Enzyme}
  (\prettyref{topic:enzym}), \I{Gerinnungsfaktoren}
  (\prettyref{topic:gerinnungsfaktoren}), Zellmembranbausteine
  u.\,a. zum Einsatz.
\end{itemize}


}{
\begin{itemize}
  \item Zuckerstoffwechsel (Kohlenhydrate)
  \begin{itemize}
    \item Energielieferant
    \item Regulation durch Insulin
  \end{itemize}
  \item Fettstoffwechsel \Z{(Fettsäuren, gesättigt und  ungesättigt)}
  \begin{itemize}
    \item Hoher Energiegehalt
    \item Lösungsmittel für Vitamine, Bausteine der Zellmembran und von
    Hormonen
  \end{itemize}
  \item Eiweißstoffwechsel (Proteine)
  \begin{itemize}
    \item Zellmembranbausteine
    \item Enzyme, Gerinnungsfaktoren, \ldots
  \end{itemize}
\end{itemize}
}{}{}{}




\cite{AlbertsMolZellbiologie2,LoefflerBiochemie7}


\section{Immunsystem}
\label{topic:immunsystem}
% TODO Inhalt: neu: Immunsystem

\Layout{0}{\TxBeschreibung}
{%
Das Immunsystem (Abwehrsystem) schützt den Körper vor
Infektionskrankheiten. Eine besonders bedeutende Rolle
spielen die \EE{Leukozyten} (Weiße Blutkörperchen,
\prettyref{topic:leukozyten}). Sie ermöglichen die Abtötung
und den Abbau von Krankheitserregern im Körper. Fieber kann
die Immunabwehr unterstützen.
}{
\begin{itemize}
  \item Schutz vor Infektionskrankheiten
  \item Leukozyten
  \item Fieber
\end{itemize}
}{}{}{}


\Layout{0}{Störungen des Immunsystems}
{%
Störungen des Immunsystems können sowohl zu einer
Überfunktion, als auch zu einer Insuffizienz führen.
Exemplarischen seien hier knochenmarkschädigende Therapien
und Allergien ausgeführt: 

\begin{itemize}
  \item \EE{Knochenmarkschädigende Therapien}:
  \index{Knochenmark!-sschädigende Therapien} Einige
  medizinische Behandlungsmethoden verursachen als Nebenwirkung eine massive \EE{Schädigung des Knochenmarks}\index{Knochenmark!Schädigung} (z.\,B. Bestrahlung, Chemotherapien bei Krebspatienten ("`onkologische Patienten"'), etc.). Diese Patienten leiden
  neben einer Blutarmut auch an einem schwachen Immunsystem, da
  nicht genügend weiße Blutkörperchen gebildet werden. Sie
  können daher sehr leicht schwerwiegende Infektionen bekommen.
  Diese Patienten
  sind im Krankentransport häufig anzutreffen.
  \item \EE{Allergie}:
  \label{topic:allergie}
  \index{Allergie}
  Das Immunsystem kann jedoch auch fälschlicherweise gegen
  an sich ungefährliche Stoffen aktiv werden.
  Man spricht dann von einer \T{Allergie}. Im leichtesten Fall
  kommt es bei Kontakt mit dem auslösenden Stoff zu einer
  Schleimhautreizung (\enquote{laufende Nase}), Augenrötung,
  tränen, einer lokalen Rötung bzw. einem Juckreiz. Je nach Stärke der
  Allergie sowie Dauer und Intensität des Kontakts kann es
  aber auch zu lebensbedrohlichen Zuständen wie dem
  anaphylaktischen Schock kommen
  (\prettyref{topic:schock-anaphylaktischer}).
\end{itemize}

}{
\begin{itemize}
  \item Knochenmarkschädigende Therapien
  \begin{itemize}
    \item Schädigung des Knochenmarks durch med. Therapien (Chemotherapien,
    Bestrahlung, \ldots)
    \item Zu wenig Leukozyten
  \end{itemize}
  \item Allergien
  \begin{itemize}
    \item Fehlgeleitete Immunreaktion gegen ungefährliche Stoffe
  \end{itemize}
\end{itemize}
}{}{}{}






\clearpage
\section{Flüssigkeit im Körper}
\index{Flüssigkeitshaushalt}
\index{Volumenhaushalt}

\Layout{2}{Verteilungsräume}
{%
Das Körperwasser macht ca. 60\PC* des Körpergewichtes
aus und teilt sich auf:

\begin{itemize}
  \item \T{IntraZellularRaum} (IZR), 40\PC* KG
  \item \T{ExtraZellularRaum} (EZR), 20\PC* KG
  \begin{itemize}
    \item \index{Zwischengewebsraum}Zwischengewebsraum,
    $\sim$\,16\PC*: "`freier"'
    Raum zwischen den Zellen im Gewebe
    \item \index{Gefäßraum}Gefäßraum $\sim$\,4\PC*:
    Raum in den Blutgefäßen
  \end{itemize}
\end{itemize}
}{}{}{}{}


\Layout{0}{Trennung}
{%
Die einzelnen Verteilungsräume haben Trennungen: Der
Intrazellularraum (\Abk{IZR}) wird vom Extrazellularraum
(\Abk{EZR}) durch die Zellmembran getrennt, der
Zwischengewebsraum vom Gefäßraum durch die Blutgefäßwand.
}{
\begin{itemize}
    \item IZR -- EZR:    Zellmembran (semipermeable Membran,
    s.\,u.)
    \item Zwischengewebsraum -- Gefäßraum:  Blutgefäßwand
\end{itemize}
}{}{}{}


\Layout{20}{Atome, Ionen und Elektrolyte}
{%
\index{Atome}%
\index{Ionen}%
\index{Elektrolyte}%
\label{topic:atome}%
\label{topic:ionen}%
\label{topic:elektrolyte}%
% 
% 
\T{Atome} sind die kleinsten Grundbestandteile von Materie.
\T{Ionen} sind Atome die \E{elektrisch geladen} sind. Sie
sind gut in Wasser löslich. Wenn sie in Wasser gelöst sind
heißen sie "`\T{Elektrolyte}"'.

Die wichtigsten Ionen sind Natrium (\ce{Na^+}), Kalium
(\ce{K^+}), Kalzium (\ce{Ca^{++}}) \Z{(positiv geladen)} und
Chlorid (\ce{Cl^-}, \Z{negativ geladen}).

\Z{Elektrolyte sind geladene Teilchen, die  bei Aufspaltung
von Säuren, Laugen oder deren Salzen in wässriger Lösung 
entstehen.  Man unterscheidet zwischen positiv geladenen
Teilchen, sog. \E{Kationen} (z.\,B. Kalium (\ce{K^+}),
Natrium (\ce{Na^+}), Magnesium (\ce{Mg^{++}}), Kalzium
(\ce{Ca^{++}}), \ldots) und negativ geladenen Teilchen, den
\E{Anionen} (z.\,B. Chlorid (\ce{Cl^-})).

Diese Ionen sind unterschiedlich verteilt: Im \Abk{IZR} ist
die \ce{K^+}-Konzentration hoch, dagegen gibt es wenig
\ce{Na^+}; Im \Abk{EZR} ist es umgekehrt: Es gibt viel
\ce{Na^+} und wenig \ce{K^+}. Spezielle Transportsysteme
sorgend dafür, dass diese Ionen derartig verteilt werden.}

\Fazit{Im Wasser gelöste, elektrisch geladene Teilchen, wie
z.\,B. Natrium, Kalium und Chlorid, heißen \T{Elektrolyte}.
Sie sind im Körper unterschiedlich verteilt.}
}{

}
{./images/wikipedia-pd/Semipermeable_membrane.png}
{Gefäßraum und Zwischengewebsraum, getrennt durch die Blutgefäßwand. Es
schwimmt viel herum: Die grünen, gelben blauen und violetten Kugeln stellen
Ionen und Moleküle dar. Die roten Scheiben sind rote Blutkörperchen.}
{WmCo/pidalka44}


\Layout{0}{Funktion der Elektrolyte}{%
Die Elektrolyte haben \E{zwei} wesentliche Funktionen:
Einerseits haben sie aufgrund der Osmose (s.\,u.) großen
Einfluss auf die \EE{Flüssigkeitsverteilung} im Körper,
andererseits haben sie durch ihre elektrische Ladung eine
entsprechende \EE{elektrische Funktion}. Durch die
elektrische Ionenladung ist jede Körperzelle gegenüber ihrer
Umgebung geladen. \Z{Würde man winzig kleine Elektroden in
eine Zelle und in ihre Umgebung stechen, wäre innen der
Minuspol und außen der Pluspol.} Dies ist besonders für die
Erregungsleitung im Herzen, für die Nerven und den Muskel
wichtig.
}{%
\begin{itemize}
    \item \EE{Verschiebung von
    Wasser} zwischen den einzelnen Räumen \Z{(durch
    Osmose und Diffusion, siehe unten)}
    \item \EE{Elektrische Funktion} durch Ionenladung. Besonders
    wichtig für
    \begin{itemize}
        \item Herz (Erregungsleitung)
        \item Nerven
        \item Muskel
    \end{itemize}
\end{itemize}
}{}{}{}





%%%%%%%%%%%%%%%%%%%%%%%%%%%%%%%%%%%%%%%%%%%%%%%%%%%%%%%%%%%
%%%%%%%%%%%%%%%%%%%%%%%%%%%%%%%%%%%%%%%%%%%%%%%%%%%%%%%%%%%
%%%%%%%%%%%%%%%%%%%%%%%%%%%%%%%%%%%%%%%%%%%%%%%%%%%%%%%%%%%
\subsection{Wasser- und Elektrolythaushalt}
\index{Wasserhaushalt}
\index{Volumenhaushalt}
\index{Elektrolythaushalt}
\label{topic:wasser-elektrolythaushalt}

\Layout{0}{Exkurs}
{%
Das Blutvolumen setzt sich zusammen aus dem \E{Plasma},
welches zum Großteil aus Wasser besteht, und den
\E{zellulären Bestandteilen} (Blutkörperchen).

Der Anteil beträgt beim Erwachsenen \VonBis{7}{8}\PC*, beim
Neugeborenes \VonBis{8}{9}\PC* des Körpergewichts.
\E{Verluste bis zu ca. 20\PC*} sind prinzipiell ohne
Funktionsbeeinträchtigung tolerierbar, allerdings kommt es
bald zu Schäden in der Endstrombahn.

Osmose und Diffusion sorgen zusammen mit den Elektrolyten
für die richtige Verteilung der Flüssigkeit in den
Verteilungsräumen und damit für die Volumenregulation.
Darüber hinaus steuern
Hormone\ZFootnote{(Renin/Angiotensin/Aldosteron)} zusammen
mit den Nieren die planmäßige Ausscheidung von Flüssigkeit.
}{
\begin{itemize}
  \item Blutvolumen = Plasma(wasser) + zelluläre Bestandteile
  \item Erwachsener \VonBis{7}{8}\PC*, Neugeborenes\,/\,Säugling \VonBis{8}{9}\PC* des Körpergewichts
  \item Verluste bis 20\PC* verkraftbar
  \item Osmose und Diffusion verteilen Wasser
  \item Hormone {\&} Nieren sorgen für Ausscheidung
\end{itemize}
}
{}{}{}









%%%%%%%%%%%%%%%%%%%%%%%%%%%%%%%%%%%%%%%%%%%%%%%%%%%%%%%%%%%
%%%%%%%%%%%%%%%%%%%%%%%%%%%%%%%%%%%%%%%%%%%%%%%%%%%%%%%%%%%
%%%%%%%%%%%%%%%%%%%%%%%%%%%%%%%%%%%%%%%%%%%%%%%%%%%%%%%%%%%
\subsection{Osmose und Diffusion}
\index{Osmose|Z}
\index{Diffusion|Z}
\label{topic:osmose}
\label{topic:diffusion}

\Definition{Verteilung von Wasser durch
Konzentrationsausgleich in verschiedene Räume, welche durch
halbdurchlässige Membranen voneinander getrennt sind.}

\AuthNotes{GAB: Osmose und Diffusion: Grundsätzlich wichtig
und relevant, allenfalls bessere Erklärung, aber bleibt in
Normalschrift, nicht heller. [GAB]

KOCH: Sollte alles grau werden}

\Layout{20}{Flüssigkeitsverschiebung durch Diffusion}
{%
%
Geladene Teilchen können die Zellmembran nicht selbstständig
durchdringen, Wasser dagegen schon. Man nennt die
Zellmembran daher eine "`\index{semipermeable
Membran}halbdurchlässige\ZFootnote{semipermeable} Membran"'.

Die Natur möchte, dass in einer Flüssigkeit überall die
gleiche Konzentration von geladenen Teilchen herrscht. Da
die Zellmembran (und teilweise auch die Blutgefäßwand) nur
das Wasser durchlässt, nicht jedoch die geladenen Teilchen,
verschiebt sich nur das Wasser von einem Raum in den
anderen.

\Z{Man es sich so vorstellen, dass die Elektrolyte einen
"`Druck"' ausüben und das Wasser dazu bringen "`die Seiten
zu wechseln"'. Man nennt diesen Druck den
"`\index{osmotischer Druck}\T{osmotischen Druck}"'. Man kann
auch sagen:
\Speech{"`Das Wasser folgt dem Konzentrationsgefälle."'}}
}{

}
% {./images/osmose.png}
{./images/wikipedia-pd/Osmotische_druk.png}
{Diffusion. }
{WmCo/Sander van der Molen, PD}





\Z{Daneben gibt es den \E{onkotischen} Druck, welcher durch
Eiweißbestandteile ausgeübt wird.~\ZFootnote{Der durch
Eiweißbestandteile hervorgerufene onkotische Druck wird bei
speziellen "`kolloidalen"' Infusionslösungen ausgenützt. Sie
enthalten große Eiweiße, welche die Blutgefäßwand
normalerweise nicht passieren können und folglich Wasser vom
Zwischengewebsraum in den Blutgefäßraum "`ziehen"'. Der
verbreitetste Wirkstoff ist die Hydroxyethylstärke (HAES),
Beispiele für damit erzeugte Infusionslösungen sind
Elohaest{\texttrademark}, Voluven{\texttrademark} und
HyperHAES{\texttrademark}.}}

\Z{Wie kommen die Elektrolyte in die einzelnen Räume wenn
doch nur Wasser die Wand passieren kann? Dafür gibt es in
der Zelle spezielle "`Pumpen"' \Z{(Transportproteine)} die
die Ionen hin und her befördern. }

Die Steuerung der richtigen Verteilung der verschiedenen
Ionen ist recht kompliziert aber lebenswichtig. Eine
"`falsche"' Verteilung der Elektrolyte kann schlimme Folgen
haben (zB. Herzrhythmusstörungen).

\Fazit{Die richtige Verteilung der Elektrolyte im Körper
ist lebenswichtig. Elektrolytstörungen können
lebensgefährlich sein. }

\Layout{0}{Häufige Ursachen von
Elektrolytstörungen}
{%
\index{Elektrolytstörungen}
\label{topic:elektrolytstoerungen}
\label{topic:dialyse-elektrolyte}
%
\begin{itemize}
    \item Niereninsuffizienz (gestörte Ausscheidung)
    \begin{itemize}
        \item Besonders bei Dialysepatienten
    \end{itemize}
    \item Durchfall, Erbrechen (Verlust von Elektrolyten)
    \item Medikamente
    \item Vergiftung
    %\item \Z{Hormonentgleisungen}
\end{itemize}
}{%
\begin{itemize}
    \item Niereninsuffizienz, Dialysepatienten
    \item Durchfall, Erbrechen (Verlust von Elektrolyten)
    \item Medikamente, Vergiftung
    %\item \Z{Hormonentgleisungen}
\end{itemize}
}{}{}{}



% \clearpage
\bigskip
%%%%%%%%%%%%%%%%%%%%%%%%%%%%%%%%%%%%%%%%%%%%%%%%%%%%%%%%%%%
%%%%%%%%%%%%%%%%%%%%%%%%%%%%%%%%%%%%%%%%%%%%%%%%%%%%%%%%%%%
%%%%%%%%%%%%%%%%%%%%%%%%%%%%%%%%%%%%%%%%%%%%%%%%%%%%%%%%%%%
\section{Säure-/Basenhaushalt}
\index{Säure-Basen-Haushalt}
\label{topic:saeure-basen-haushalt}

\subsection{Allgemeines}

\Layout{abstract}{}
{Der Körper muss einen ausgewogenen Säure-Basenhaushalt
(pH-Wert) haben. Puffersysteme sind wichtige Intrumente zur
Regulation. Der pH-Wert hängt eng mit der Atmung zusammen:
Säure kann zu Wasser und \ce{CO2} abgebaut werden. Bei zu
geringer Atmung sammelt sich im Körper Säure, bei
viel Atmung wird mehr Säure verbraucht. Umgekehrt erhöht
sich bei Übersäuerung der Atemantrieb. 

}{}{}{}{}

\Layout{0}{Säuren und Laugen}
{%
\index{Säure}
\index{Lauge}
\index{Base}
\T{Säuren} sind Stoffe, die Protonen\footnote{\E{Proton}:
positiv geladenes Teilchen im Atoms. Ein positiv geladenes
Wasserstoffion (\ce{H^+}) entspricht der Ladung eines Protons, daher kann man vereinfacht sagen:
\Speech{"`\ce{H^+} entspricht Säure."'}} abgeben können,
\T{Basen} sind Stoffe, die Protonen aufnehmen können. Der
Begriff
\T{Lauge} ist gleichbedeutend mit dem Begriff Base. \cite{LoefflerBiochemie7}

\bigskip

\bigskip
}{%
\begin{description}
  \item[\T{Säuren}]: Stoffe, die
  Protonen (\ce{H^+}) abgeben können
  \item[\T{Basen}\,/\,\T{Laugen}]: Stoffe, die
  Protonen (\ce{H^+}) aufnehmen können
\end{description}
}{}{}{}



\Layout{0}{pH-Wert}
{%
\index{pH-Wert}
Der pH-Wert gibt an, wie \E{sauer} oder \E{basisch} eine
Lösung ist. Ein pH-Wert {\textless}\,7 gilt als \E{sauer},
ein pH\,=\,7 als \E{neutral} und ein ph\,{\textgreater}\,7
als \E{basisch}.

Der Körper ist da
sehr empfindlich, der pH-Wert muss sich in engen Grenzen
bewegen. Schon geringe Abweichungen können große Folgen
haben. Bei einer Übersäuerung sprechen wir von einer
\T{Azidose}, wenn der Patient basisch wird,
von einer \T{Alkalose}.
 
% \Fazit{Bei einer \EE{Übersäuerung} sprechen wir von einer
% \T{Azidose}, wenn der Patient \EE{basisch} wird von einer
% \T{Alkalose}.}

\Z{Der richtige pH-Wert ist lebenswichtig. Er wird beim Menschen
normalerweise vom Körper sehr konstant um ca 7,4 gehalten. Bereits einen pH-Wert
unter 7,35 (!)  bezeichnet man als
\E{Azidose}, bzw. über 7,45  als \E{Alkalose}. Leider gibt
es präklinisch normalerweise keine Möglichkeit den pH-Wert zu messen, jedoch hat die Atmung
bzw. Beatmung sowie Erkrankungen
 großen Einfluss auf ihn, siehe unten.} 
}{
\begin{itemize}
  \item Gibt an, wie sauer etwas ist.
  \item pH\,{\textless}\,7 {\dotfill} sauer, 
  pH\,=\,7 ist neutral,
  \item pH\,{\textgreater}\,7 {\dotfill} basisch
\end{itemize}
}{}{}{}

\Layout{0}{Regulationsmechanismen}
{%
Der Körper muss den pH-Wert in einem engen Bereich halten.
Dazu gibt es Regulationsmechanismen. Der wichtigste
Regulationsmechanismus ist das
\Z{Bicarbonat-}\II{Puffersystem} im Blut, an
dem auch die \EE{Atmung} und die \EE{Lunge} wesentlich
beteiligt sind.\ZFootnote{Ein Puffer ist ein Stoff welcher
pH-Änderungen abfedert.} Daneben gibt es auch noch andere
Regulationsmechanismen, so reguliert auch die \EE{Niere} die
Ausscheidung oder Wiederaufnhame von Säure in Form
von Wasserstoffionen.

}{
\begin{itemize}
  \item Bikarbonat-Puffer\begin{itemize}
    \item Blut, Lunge (Atmung)
  \end{itemize}
  \item Niere (\ce{H^+-Ausscheidung})
\end{itemize}
}{}{}{}

\Layout{0}{Puffersystem im Blut}
{%
\label{topic:bikarbonatpuffer}
Der  \Z{Bicarbonat-}Puffer wandelt Säure  und
\Z{Bikarbonat} in Kohlendioxid \Z{\ce{CO2}} und
Wasser \Z{\ce{H2O}} um. Dieser Mechanismus funktioniert in
beide Richtungen, je nachdem welcher Stoff (Säure bzw.
\ce{CO2}) vermehrt vorliegt.
\Commentary[Säure-Basen-Haushalt / Bikarbonat-Puffer /
Formel]{Die vollständige Formel lautet:

\protect\ce{HCO3^- + H^+ <->[\text{Gleichgewicht}] CO2 + H2O}

Dies ist die konkrete Ableitung der
Henderson-Hasselbach'schen Puffergleichung für den
Bikarbonatpuffer.
}

\bigskip
\Formula{\centering
% \begin{displaymath}
\ce{$\text{\Z{\tiny Bikarbonat + }}$ $\text{Säure}$
        <->[\text{Gleichgewicht}] CO2 $\text{\Z{\tiny + Wasser}}$}
% \end{displaymath}
}
        
Das \ce{CO2} wird normalerweise über die Lunge abgeatmet. Daraus folgt: 

\begin{itemize}
  \item Atmet der Patient nicht (genug), \E{staut}
  sich das \ce{CO2} und Säure wird auf der anderen Seite \E{nicht
  abgebaut}. Es entsteht eine \EE{atmungsbedingte Übersäuerung}.
  \item Atmet der Patient zu viel, wird zu viel
  \ce{CO2} abgeatmet und wird nachgebildet; dabei wird Säure
  zu Wasser und geht "`verloren"'. Da nun Säure 
  \E{fehlt}, kommt es zu einer \EE{atmungsbedingten Alkalose}.
  \item Fällt zu viel Säure im
  Körper an (Stoffwechsel-bedingte Übersäuerung), wird
  es zu  \ce{CO2} (und Wasser) umgewandelt, \E{das \ce{CO2} wird
  abgeatmet.} Es kommt zu einer \EE{tiefen, schnellen Atmung}.
\end{itemize}
}{%
\ce{$\text{\Z{\tiny Bik. + }}$ $\text{Säure}$
  <->[\text{Gleichgewicht}] CO2 $\text{\Z{\tiny + Wasser}}$  }

\ce{CO2} wird von Lunge abgeatmet.
\begin{itemize}
  \item Zu wenig Atmung: \ce{CO2} staut
  sich, \ce{H^+} wird nicht abgebaut
  
  \EE{Atmungsbedingte Übersäuerung}
  \item Zu viel Atmung:  zu viel \ce{CO2} abgeatmet, \ce{CO2} wird
  nachgebildet, Säure wird zu Wasser und fehlt
  
  \EE{Atmungsbedingte Alkalose}
  \item Zu viel Säure  im
  Körper: zu \ce{CO2} umgewandelt und wird abgeatmet
\end{itemize}
}{}{}{}








% \clearpage % TEMP temp clearpage
\bigskip
\subsection{Auswirkungen der Atmung und des Stoffwechsels
auf den Säure-Basen-Haushalt}


\Layout{0}{Atmungsbedingte Übersäuerung (Respiratorische
Azidose)} {%
\label{bkm:RefHeading35233321}%
\label{topic:azidose-respiratorische}%
\index{Uebersäuerung@Übersäuerung}%
\index{Respiratorische Azidose}%
Zugrunde liegt eine
insuffiziente innere oder äußere Atmung. Dies Umfasst z.\,B.  eine zu
niedrige Atemfrequenz  oder ein zu geringes
Atemzugsvolumen, eine Verlegung der Atemwege,
ein zu geringes Sauerstoffangebot in der Einatemluft, einen
beeinträchtigten Gasaustausch oder generell eine
Kreislaufinsuffizienz.

Da die \ce{CO2}-Abatmung vermindert ist, kommt es zu einer
Anhäufung von \ce{CO2} im Blut \Z{(Hyperkapnie)} und
folglich zu einer Ansäuerung des Blutes.

\bigskip

\bigskip

\bigskip

\bigskip
}
{\begin{itemize}
  \item Insuffiziente innere oder äußere Atmung
  \begin{itemize}
    \item AF $\downarrow$, AZV $\downarrow$
    \item Sauerstoffgehalt der Einatemluft $\downarrow$
    \item Gasaustausch beeinträchtigt
    \item Kreislaufinsuffizienz
  \end{itemize}
  \item Anhäufung von
\ce{CO2} im Blut
\end{itemize}}
{}{}{}

% \MassnahmenNG{Spezielle Maßnahmen}{Respiratorische Azidose}
% {}{
% \index{Alkalose!Respiratorische}
% 
% \begin{itemize}
%     \item Behebung der Ursache
%     \item Erhöhung des Atem-Minuten-Volumens
% \end{itemize}
% }







\Layout{0}{Stoffwechselbedingte Übersäuerung
(Metabolische Azidose)}
{%
\label{topic:azidose-metabolische}
\index{Uebersäuerung@Übersäuerung}
\index{Azidose!Metabolische}
\index{Hyperventilation!bei metabolischer Azidose}
Es kommt zu einer übermäßigen Säureproduktion des Körpers.
Ein klassisches Krankheitsbild welches dazu führt  ist das
\II[Koma!diabetisches]{Diabetische Koma}. Hier werden
saure Stoffe gebildet \Z{(sog. Ketonkörper)}, welche ein
Ungleichgewicht verursachen. Es kann dadurch zu einer Schädigung des Zellstoffwechsels
und des Herzmuskels kommen. 

Es bestehen die Symptome der Grunderkrankung. Weiters fällt
eine \EE{tiefe, schnelle Atmung} auf, da der Körper versucht,
\ce{CO2} abzuatmen (\E{Kußmaul'sche Atmung}). 



\bigskip

\bigskip

}{%
\begin{itemize}
  \item Bildung von sauren Stoffen im Körper
  
  Z.\,B. diabetisches Koma
  \item Symptome der Grunderkrankung
  \item Tiefe, schnelle Atmung
  \item Schädigung des Zellstoffwechsels und des
  Herzmuskels
\end{itemize}
}{}{}{}


\Layout{2}
{Atmungsbedingte (respiratorische) Alkalose}
{\label{topic:alkalose-respiratorische}
\index{Alkalose}
\index{Hyperventilation!respiratorische Alkalose}

Durch die vermehrte Abatmung von \ce{CO2} wird der pH-Wert
des Blutes basisch und es kommt zu einem
Elektrolytungleichgewicht. Dies verursacht Symtome wie
Bewusstseinsstörungen, Kribbeln in den Fingern,
Pfötchenstellung und tonischen Krämpfe.

Eine ausführliche Beschreibung des
\E{Hyperventilationssyndroms} findet sich unter
\prettyref{topic:hyperventilationstetanie}.}
{}
{}
{}
{}

\Layout{2}
{Stoffwechselbedingte (metabolische) Alkalose}
{\label{topic:alkalose-metabolische}
\index{Alkalose!Metabolische}
\index{Alkalose!Stoffwechselbedingte}



Eine stoffwechselbedingte Alkalose  kommt v.a. durch den \E{Verlust von Säure},
z.\,B. durch anhaltendes Erbrechen, zu Stande.}
{}
{}
{}
{}


\subsubsection*{Zusammenfassung}


\Fazit{Die \EE{richtige Verteilung} von Wasser im Körper ist
lebenswichtig.}

\Fazit{Mittels Osmose und Diffusion verteilt der Körper das
Wasser.} 

\Fazit{Die \EE{Elektrolyte} haben einen großen Einfluss auf die
  Wasserverteilung: Nur wenn die Elektrolyte richtig verteilt sind, kann das
  Wasser richtig verteilt werden.}
\Fazit{Elektrolyte haben auch
  eine elektrische Funktion.}
\Fazit{Ein ausgeglichener Säure-Basen-Haushalt ist wichtig, damit chemische
  Abläufe im Körper richtig funktionieren können.}
  
\Fazit{Besonders die \EE{Atmung} und der \EE{Stoffwechsel} können dem
  Säure-Basen-Haushalt Streiche spielen: \newline
  Wer zu wenig atmet wird sauer. \newline
  Wer zu viel atmet wird basisch und spürt ein Kribbeln
in den Fingern. \newline
Wer Probleme mit dem Stoffwechsel hat kann auch
    sauer oder basisch werden.}







% %%%%%%%%%%%%%%%%%%%%%%%%%%%%%%%%%%%%%%%%%%%%%%%%%%%%%%%%%%%
% %%%%%%%%%%%%%%%%%%%%%%%%%%%%%%%%%%%%%%%%%%%%%%%%%%%%%%%%%%%
% %%%%%%%%%%%%%%%%%%%%%%%%%%%%%%%%%%%%%%%%%%%%%%%%%%%%%%%%%%%
% %%%%%%%%%%%%%%%%%%%%%%%%%%%%%%%%%%%%%%%%%%%%%%%%%%%%%%%%%%%
% \section{Tod}
% 
% 
% 
% Der Tod ist eine der grundlegensten und unausweichlichsten
% Vitalfunktionen überhaupt. Vgl.   \prettyref{chp:TOD}.

\nocite{AMLS3,emergency9,PHTLS6,RedelsteinerHRNS1}

\ChapterBibliography




